% 来源：SGA 4-3 重排本 PDF 第 1--37 页；第 IX 讲。
\SourcePage{1}
\begin{center}
M. Artin、A. Grothendieck、J.-L. Verdier\\
N. Bourbaki、P. Deligne、B. Saint-Donat 参与\\[1em]
\textbf{\Term{topos}{拓扑斯}理论与\Term{scheme}{概形}的\Term{etale-cohomology}{平展上同调}}\\
第一册\\
\Term{topos}{拓扑斯}理论
\end{center}
\SourcePage{2}
M. Artin、A. Grothendieck、J.-L. Verdier；N. Bourbaki、P. Deligne、B. Saint-Donat 参与。
\SourcePage{3}
\begin{center}
\textbf{\Term{topos}{拓扑斯}理论与\Term{scheme}{概形}的\Term{etale-cohomology}{平展上同调}}\\
第一册：\Term{topos}{拓扑斯}理论\\
M. Artin、A. Grothendieck、J.-L. Verdier\\
N. Bourbaki、P. Deligne、B. Saint-Donat 参与
\end{center}
\textbf{摘要。}——

\textbf{英文摘要。}——待补。
\SourcePage{4}
% 原页空白。
\SourcePage{5}
% 原总目录第 5--7 页由译本目录重建。
\SourcePage{6}
\SourcePage{7}
\SourcePage{8}
\SGAchapter{IX}{可构造层；代数曲线的上同调}{Faisceaux constructibles. Cohomologie d'une courbe algébrique}
\begin{center}M. Artin\end{center}
% 本讲目录由译本目录重建。
\SGAsection{0}{引言}
应当注意，从本讲开始，与前面各讲以及\Term{topological-space}{拓扑空间}的经典\Term{cohomology}{上同调}理论不同，为使主要命题成立，我们不得不把范围\Term{restriction}{限制}于\Term{torsion-sheaf}{挠层}（而且常常更精确地\Term{restriction}{限制}于其\Term{torsion}{挠性}与各\Term{residue-characteristic}{剩余特征}互素的\Term{torsion-sheaf}{层}）。

指出一点：实际上，对于\Term{characteristic}{特征} $p\ne0$ 的\Term{algebraically-closed-field}{代数闭域}上的\Term{algebraic-variety}{代数簇}，并不存在一种与 $k=\mathbb C$ 时的经典理论类似的、以整数（甚至实数）为系数的“合理”\Term{cohomology}{上同调}理论（Serre 已经指出这一点）。更精确地说，不存在这样一个\Term{contravariant-functor}{反变函子} $H^1$：例如定义在\Term{characteristic}{特征} $p>0$ 的\Term{algebraically-closed-field}{代数闭域} $K$ 上的\Term{smooth}{光滑}\Term{projective}{射影}\Term{scheme}{概形}\Term{category}{范畴}上，取值于 $\mathbb R$（或 $\mathbb R$ 的一个子域）上的\Term{finite}{有限}维\Term{vector-space}{向量空间}\Term{category}{范畴}，“与积交换”，即满足
\[
H^1(X\times Y)\xleftarrow{\sim}H^1(X)\times H^1(Y),
\]
并且当 $X$ 是\Term{genus}{亏格}为 $1$ 的\Term{connected}{连通}曲线时，$\dim H^1(X)=2$。事实上，由此将推出：若 $X$ 是 $k$ 上的\Term{abelian-variety}{阿贝尔簇}，则从 $\operatorname{End}(X)$ 到 $\operatorname{End}(H^1(X))$ 的映射 $u\mapsto u^1$ 是可加的，因而给出 $\operatorname{End}(X)$ 的反环的一个表示。
\SourcePage{9}
然而，在\Term{characteristic}{特征} $p$ 时，存在\Term{elliptic-curve}{椭圆曲线} $X$，其\Term{endomorphism}{自同态}环是在 $\mathbb Z$ 上定义的某个四元数\Term{algebra}{代数}的一个\Term{maximal-order}{极大序} $A$（[3]，第 198 页）；这种\Term{algebra}{代数}显然不可能在 $\mathbb R$ 上的二维\Term{vector-space}{向量空间}中有表示，因为 $A\otimes_{\mathbb Z}\mathbb R$ 是一个除环（四元数体）。

\SGAsection{1}{挠层的若干基本性质}
设 $T$ 为一个\Term{topos}{拓扑斯}，$F$ 为 $T$ 上的一个\Term{abelian-sheaf}{阿贝尔层}。可在 $F$ 中作乘以 $n$ 的运算（$n\in\mathbb Z$ 为整数），它由 $(\mathrm{Ab})$ 中的相应运算诱导。用 ${}_nF$ 表示这一乘法的核；因此，对于 $X\in\operatorname{Ob}T$，${}_nF(X)$ 是 $F(X)$ 中阶整除 $n$ 的\Term{section}{截面}所成的群。

用 $\mathbb P$ 表示所有素数的\Term{set}{集合}。
\begin{SGAresult}{定义}{1.1}
设 $\mathfrak p$ 为一个素数\Term{set}{集合}。若\Term{canonical}{典范}\Term{morphism}{态射}
\[
\varinjlim_n{}_nF\longrightarrow F
\]
是\Term{bijection}{双射}，就称 $F$ 是\Term{p-torsion-sheaf}{$\mathfrak p$-挠层}，或称为\Term{p-torsion-sheaf}{$\mathfrak p$-层}；这里 $n$ 遍历满足 $\operatorname{ass}n\subset\mathfrak p$ 的整数，即 $n$ 的素因子全属于 $\mathfrak p$。若 $\mathfrak p=\mathbb P$ 是所有素数的\Term{set}{集合}，就简称 $F$ 为\Term{torsion-sheaf}{挠层}。
\end{SGAresult}
\begin{SGAresult}{命题}{1.2}
(i) $F$ 是\Term{p-torsion-sheaf}{$\mathfrak p$-挠层}，当且仅当 $F$ 是一个以 $\mathfrak p$-挠\Term{abelian-group}{阿贝尔群}为值的\Term{presheaf}{预层}的\Term{adjunction}{伴随}\Term{sheaf}{层}；可取
\[
P=\underset{\operatorname{ass}n\subset\mathfrak p}{\varinjlim{}^{(\mathrm{pr\text{é}f})}}{}_nF.
\]
(ii) 若 $F$ 是\Term{p-torsion-sheaf}{$\mathfrak p$-挠层}，而 $X\in\operatorname{Ob}T$ 是\Term{quasi-compact}{拟紧}\Term{object}{对象}\footnote{若每个\Term{covering-family}{覆盖族} $\{X_i\to X\}$ 都可由一个\Term{finite-covering-family}{有限覆盖族}\Term{refinement}{加细}，就称\Term{topos}{拓扑斯}中的\Term{object}{对象} $X$ 为\Term{quasi-compact}{拟紧}的。若每个\Term{quasi-compact}{拟紧}\Term{object}{对象}的\Term{inverse-image}{逆像}仍然\Term{quasi-compact}{拟紧}，就称\Term{topos}{拓扑斯}\Term{morphism}{态射} $u$ 为\Term{quasi-compact}{拟紧}的。}，则 $F(X)$ 是一个 $\mathfrak p$-\Term{torsion-group}{挠群}。

(iii) 若 $T$ \Term{local-scheme}{局部}为\Term{finite-type}{有限型}（第 VI 讲 1.1），则 $F$ 是\Term{p-torsion-sheaf}{$\mathfrak p$-挠层}，当且仅当对每个\Term{quasi-compact}{拟紧}的 $X\in\operatorname{Ob}T$，$F(X)$ 都是 $\mathfrak p$-\Term{torsion-group}{挠群}。在此情形，对每个\Term{quasi-compact}{拟紧}的 $X$ 及所有 $q$，有
\[
\underset{\operatorname{ass}n\subset\mathfrak p}{\varinjlim}H^q(T/X;{}_nF)
 \xrightarrow{\sim}H^q(T/X;F).
\]
(iv) 若 $u:T\to T'$ 是\Term{topos}{拓扑斯}\Term{morphism}{态射}，而 $F'$ 是 $T'$ 上的\Term{p-torsion-sheaf}{$\mathfrak p$-挠层}，则\Term{inverse-image}{逆像} $u^*F$ 是 $T$ 上的\Term{p-torsion-sheaf}{$\mathfrak p$-挠层}。
\SourcePage{10}
(v) 若 $u:T\to T'$ 是\Term{topos}{拓扑斯}\Term{morphism}{态射}，$T,T'$ 都\Term{local-scheme}{局部}为\Term{finite-type}{有限型}，$u$ 为\Term{quasi-compact}{拟紧}\Term{morphism}{态射}，而 $F$ 是 $T$ 上的\Term{p-torsion-sheaf}{$\mathfrak p$-挠层}，则对所有 $q$，$R^qu_*F'$ 都是 $T$ 上的\Term{p-torsion-sheaf}{$\mathfrak p$-挠层}。
\end{SGAresult}
\begin{proof}
(i) 显然，若\Term{sheaf}{层} $F$ 是\Term{p-torsion-sheaf}{$\mathfrak p$-挠层}，则 $F$ 是下述\Term{presheaf}{预层} $P$ 的\Term{adjunction}{伴随}\Term{sheaf}{层}：其\Term{section}{截面}群 $P(X)$ 是 $F(X)$ 中所有 $\mathfrak p$-\Term{torsion-element}{挠元素}所成的\Term{subgroup}{子群}。反过来，设 $F=aP$（第 II 讲）是 $P$ 的\Term{adjunction}{伴随}\Term{sheaf}{层}，其中 $P$ 取值于 $\mathfrak p$-挠\Term{abelian-group}{阿贝尔群}\Term{category}{范畴}。由\Term{exact-sequence}{正合列}
\[
0\longrightarrow{}_nP\longrightarrow P\xrightarrow{n}P
\]
得到\Term{exact-sequence}{正合列}
\[
0\longrightarrow a({}_nP)\longrightarrow F\xrightarrow{n}F,
\]
故 $a({}_nP)={}_nF$。因为\Term{functor}{函子} $a$ 与\Term{inductive-limit}{归纳极限}交换，由
\[
\underset{\operatorname{ass}n\subset\mathfrak p}{\varinjlim{}^{(\mathrm{pr\text{é}f})}}{}_nP\xrightarrow{\sim}P
\]
可知也有
\[
\underset{\operatorname{ass}n\subset\mathfrak p}{\varinjlim}{}_nF\xrightarrow{\sim}F.
\]
(ii) 置 $P=\varinjlim{}^{(\mathrm{pr\text{é}f})}{}_nF$。则 $P\subset F$，所以 $P$ 是\Term{separated}{分离}\Term{presheaf}{预层}，因而
\[
aP(X)=\varinjlim\left(\ker\left(\prod_iP(X_i)\rightrightarrows\ \right)\right),
\]
这里同通常一样，\Term{limit}{极限}沿\Term{covering-family}{覆盖族} $\{X_i\to X\}$ 取。由于 $X$ \Term{quasi-compact}{拟紧}，只需取\Term{finite-covering-family}{有限覆盖族}。对于这些族，$\prod_iP(X_i)$ 是 $\mathfrak p$-\Term{torsion-group}{挠群}，因此核也是如此。故 $aP(X)=F(X)$ 是 $\mathfrak p$-\Term{torsion-group}{挠群}。

(iii) 为验证 $\varinjlim{}_nF\to F$ 为\Term{bijection}{双射}，只须对某一\Term{family-of-generators}{生成元族}中的每个 $X\in\operatorname{Ob}T$，在 $X$ 上取值后验证。因此第一条断言归结为 (ii)。第二条断言是下述结果的推论：\SGAsourcegap{原重排本此处缺少引用；补源原版第 IX 讲第 4 页同样留空。}

(iv) 由 (i) 得到，因为 $u^*(aP')=a(u\mathbin{\cdot}P)$。

(v) 容易由 (iii) 得到。
\end{proof}
还将需要一个非阿贝尔变体。不过其定义有一些技术问题（我不知道它们是否严重），因此这里只对一个\Term{scheme}{概形}的\Term{etale}{平展}\Term{topology}{拓扑}给出定义。
\begin{SGAresult}{定义}{1.3}
设 $\mathfrak p$ 是一个素数\Term{set}{集合}。若 $G$ 的每个\Term{finite}{有限}\Term{subset}{子集}都生成一个阶为 $n$、满足 $\operatorname{ass}n\subset\mathfrak p$ 的\Term{finite}{有限}\Term{subgroup}{子群}，就称群 $G$ 为\Term{ind-p-group}{ind-$\mathfrak p$-群}。这\Term{equivalence}{等价}于说，阶为 $n_\alpha$ 且满足 $\operatorname{ass}n_\alpha\subset\mathfrak p$ 的\Term{finite}{有限}\Term{subgroup}{子群} $G_\alpha$ 构成一个\Term{filtered}{滤过}集，而且 $G=\varinjlim G_\alpha$。若 $\mathfrak p=\mathbb P$，则以\Term{ind-finite-group}{ind-有限群}代替\Term{ind-p-group}{ind-$\mathfrak p$-群}这一称呼。
\end{SGAresult}
\SourcePage{11}
立即可以验证：
\begin{SGAresult}{命题}{1.4}
(i) \Term{ind-p-group}{ind-$\mathfrak p$-群}的\Term{subgroup}{子群}及\Term{quotient-group}{商群}仍然是\Term{ind-p-group}{ind-$\mathfrak p$-群}。

(ii) \Term{ind-p-group}{ind-$\mathfrak p$-群}的\Term{filtered}{滤过}\Term{inductive-limit}{归纳极限}是\Term{ind-p-group}{ind-$\mathfrak p$-群}。

(iii) \Term{ind-p-group}{ind-$\mathfrak p$-群}的\Term{finite}{有限}\Term{projective-limit}{投射极限}是\Term{ind-p-group}{ind-$\mathfrak p$-群}。
\end{SGAresult}
\begin{SGAresult}{定义}{1.5}
设 $X$ 为一个\Term{scheme}{概形}。若对每个\Term{etale}{平展}\Term{morphism}{态射} $U\to X$，其中 $U$ \Term{quasi-compact}{拟紧}，$F(U)$ 都是\Term{ind-p-group}{ind-$\mathfrak p$-群}（分别为\Term{ind-finite-group}{ind-有限群}），就称 $X$ 上的群\Term{sheaf}{层} $F$ 是\Term{ind-p-group}{ind-$\mathfrak p$-群}\Term{sheaf}{层}（分别为\Term{ind-finite-group}{ind-有限群}\Term{sheaf}{层}）。
\end{SGAresult}
\begin{SGAresult}{命题}{1.6}
设 $F$ 是\Term{scheme}{概形} $X$ 上的一个群\Term{sheaf}{层}。

(i) $F$ 是\Term{ind-p-group}{ind-$\mathfrak p$-群}\Term{sheaf}{层}，当且仅当对 $X$ 的每个\Term{geometric-point}{几何点} $\xi$，\Term{fiber}{纤维} $F_\xi$ 都是\Term{ind-p-group}{ind-$\mathfrak p$-群}。

(ii) 若 $f:X\to X'$ 是\Term{morphism}{态射}，而 $F'$ 是 $X'$ 上的\Term{ind-p-group}{ind-$\mathfrak p$-群}\Term{sheaf}{层}，则 $f^*F'$ 是\Term{ind-p-group}{ind-$\mathfrak p$-群}\Term{sheaf}{层}。

(iii) 若 $f:X\to X'$ 是\Term{quasi-compact}{拟紧}\Term{morphism}{态射}，而 $F$ 是 $X$ 上的\Term{ind-p-group}{ind-$\mathfrak p$-群}\Term{sheaf}{层}，则 $f_*F$ 是\Term{ind-p-group}{ind-$\mathfrak p$-群}\Term{sheaf}{层}。
\end{SGAresult}
\begin{proof}
(i) 设 $F$ 是\Term{ind-p-group}{ind-$\mathfrak p$-群}\Term{sheaf}{层}，$\xi$ 是 $X$ 的\Term{geometric-point}{几何点}。则 $F_\xi=\varinjlim F(X')$，其中 $X'$ 遍历 $X$ 上的一个\Term{etale}{平展}\Term{scheme}{概形}的\Term{pseudo-filtered-category}{伪滤过}系统（第 VIII 讲 3.9）。显然，其中\Term{affine}{仿射}的（因而\Term{quasi-compact}{拟紧}的）$X'$ 构成一个\Term{cofinal}{余尾}集，因此 $F_\xi$ 是\Term{ind-p-group}{ind-$\mathfrak p$-群}的\Term{inductive-limit}{归纳极限}，由 1.4 (ii)，它是\Term{ind-p-group}{ind-$\mathfrak p$-群}。反过来，设对每个 $\xi$，\Term{fiber}{纤维} $F_\xi$ 都是\Term{ind-p-group}{ind-$\mathfrak p$-群}；又设 $U\to X$ \Term{etale}{平展}，$U$ \Term{quasi-compact}{拟紧}。取\Term{finite}{有限}\Term{subset}{子集} $S\subset F(U)$。对 $U$ 的每个\Term{geometric-point}{几何点} $\xi$，$S$ 在 $F_\xi$ 中的像生成一个\Term{finite-group}{有限群}。由于\Term{finite-group}{有限群}为\Term{finite-presentation}{有限表现}，容易推出存在 $U$ 上带 $\xi$ 标点的\Term{etale}{平展}\Term{object}{对象} $U_\xi$，使 $S$ 在 $F(U_\xi)$ 中生成一个\Term{finite-group}{有限群}（参看第 VIII 讲 4）。\Term{finite}{有限}个这样的 $U_\xi$ 即构成 $U$ 的一个\Term{covering}{覆盖}，记为 $\{U_1,\ldots,U_n\}$。则 $S$ 在 $\prod_iF(U_i)$ 中的像生成一个\Term{finite}{有限}\Term{subgroup}{子群}；因为 $F(U)\subset\prod_iF(U_i)$，结论得证。

断言 (ii) 由 (i) 立即得到，(iii) 也是直接的。
\end{proof}
\SGAsection{2}{可构造层}
\textbf{2.0.} 设 $T$ 为一个\Term{topos}{拓扑斯}。回忆：对每个 $S\in\operatorname{Ob}(\mathrm{Ens})$，有\Term{adjunction}{伴随}于下述\Term{presheaf}{预层} $P$ 的\Term{sheaf}{层} $S_T$：对每个 $X\in\operatorname{Ob}(T)$，$P(X)=S$。表示 $S_T$ 的 $T$ 中\Term{object}{对象}为
\[
\coprod_{s\in S}e=\text{“}S\times e\text{”},
\]
其中 $e$ 是 $T$ 的\Term{final-object}{终对象}。称 $S_T$ 为值为 $S$ 的\Term{constant-sheaf}{常层}（第 IV 讲）。同样可以自然地定义值为 $M$ 的\Term{constant-group}{常群}\Term{sheaf}{层}、常 $A$-模\Term{sheaf}{层}（$A$ 为环）和\Term{constant-sheaf}{常阿贝尔层}。
\SourcePage{12}
若\Term{constant-sheaf}{常层}之间的\Term{morphism}{态射} $f_T:S_T\to S'_T$ 来自一个\Term{set}{集合}\Term{morphism}{态射} $f:S\to S'$，就称它为常\Term{morphism}{态射}。

若存在\Term{final-object}{终对象}的一个\Term{covering}{覆盖} $\{e_i\to e\}$，使 $F$ 在每个 $e_i$ 上变为\Term{constant-sheaf}{常层}，就称\Term{sheaf}{层} $F$ 为\Term{locally-constant}{局部常值}的。类似地定义\Term{locally-constant}{局部常值}的群\Term{sheaf}{层}或 $A$-模\Term{sheaf}{层}，以及在 $F,G$ 为\Term{locally-constant}{局部常值}\Term{sheaf}{层}时的\Term{locally-constant}{局部常值}\Term{morphism}{态射} $f:F\to G$。最后，若\Term{locally-constant}{局部常值}的群\Term{sheaf}{层}或 $A$-模\Term{sheaf}{层}的\Term{local-scheme}{局部}取值为\Term{finite-type}{有限型}（分别为\Term{finite-presentation}{有限表现}）群或 $A$-模，就称此\Term{sheaf}{层}为\Term{finite-type}{有限型}（分别为\Term{finite-presentation}{有限表现}）的。
\begin{SGAresult}{引理}{2.1}
(i) 设 $f:F\to G$ 是 $T$ 上\Term{locally-constant}{局部常值}的\Term{set}{集合}\Term{sheaf}{层}之间的\Term{morphism}{态射}，且 $F$ 的\Term{local-scheme}{局部}取值为\Term{finite}{有限}\Term{set}{集合}。则 $f$ 是\Term{locally-constant}{局部常值}的。

(ii) 设 $f:F\to G$ 是\Term{locally-constant}{局部常值}的 $A$-模\Term{sheaf}{层}之间的\Term{morphism}{态射}，且 $F$ 为\Term{finite-type}{有限型}。则 $f$ 是\Term{locally-constant}{局部常值}的，其核与\Term{cokernel}{余核}都是\Term{locally-constant}{局部常值}的。

(iii) 设 $0\to F'\to F\to F''\to0$ 是 $A$-模\Term{sheaf}{层}（分别为群\Term{sheaf}{层}）的\Term{exact-sequence}{正合列}，$F',F''$ \Term{locally-constant}{局部常值}，$F''$ 为\Term{finite-presentation}{有限表现}。则 $F$ \Term{locally-constant}{局部常值}。
\end{SGAresult}
\begin{proof}[{
(i) 和 (ii) 的证明}]
容易归结到 $F,G$ 都是\Term{constant-sheaf}{常层}的情形。处理 (ii)：设 $F=M_T$，$G=N_T$（$M,N$ 为 $A$-模），$S\subset M$ 为 $M$ 的一个\Term{finitely-generated}{有限生成}集。则 $f$ 由其在 $S_T$ 上的\Term{restriction}{限制}决定；若 $f|S_T$ \Term{locally-constant}{局部常值}，则 $f$ 也\Term{locally-constant}{局部常值}。因此，(ii) 的第一条断言归结为 (i)，第二条断言则由第一条立即得到。

所以只需证明 (i)。设 $F=S_T$、$G=S'_T$ 都为\Term{constant-sheaf}{常层}，仍以 $f$ 记\Term{object}{对象}之间的\Term{morphism}{态射} $f:S\times e\to S'\times e$（\Term{signature}{符号}如上）。记 $f_s:e\to S'\times e$（$s\in S$）为 $f$ 的各\Term{set-component}{分量}：由于 $S$ \Term{finite}{有限}，显然只需处理 $f=f_s$ 的情形，即 $S$ 是单元素\Term{set}{集合}。设 $X_{s'}\to e$（$s'\in S'$）为使下图为\Term{cartesian-diagram}{笛卡儿图}的\Term{morphism}{态射}族：
\[
\begin{tikzcd}
X_{s'}\arrow[r]\arrow[d]&e\arrow[d,"i_{s'}"]\\
e\arrow[r,"f"]&S'\times e.
\end{tikzcd}
\]
其中 $i_{s'}:e\to S'\times e$ 是到第 $s'$ 个\Term{set-component}{分量}的“包含”。族 $\{i_{s'}\}$ 显然是\Term{covering-family}{覆盖族}，所以 $\{X_{s'}\to e\}$ 也是。由于\Term{topos}{拓扑斯}中的\Term{direct-sum}{直和}互\Term{disjoint}{不相交}，可立即验证 $f$ 在 $X_{s'}$ 上变成常\Term{morphism}{态射}，因此 $f$ 是\Term{locally-constant}{局部常值}的。
\end{proof}
\SourcePage{13}
\begin{proof}[{(iii) 的证明}]
处理 $A$-模\Term{sheaf}{层}的情形。可设 $F',F''$ 都为\Term{constant-sheaf}{常层}，$F''$ 的值为\Term{finite-presentation}{有限表现}的 $A$-模 $M''$。若 $M''$ 是\Term{free-module}{自由模}，因而有一个\Term{finite}{有限}基 $(e_i)_{1\le i\le n}$，则 $e_i$ 定义 $F''$ 的\Term{section}{截面}，因而\Term{local-scheme}{局部}可\Term{lifting}{提升}为 $F$ 的\Term{section}{截面}。这表明 $F''$ 被 $F'$ 的\Term{group-extension}{扩张} $F$ \Term{local-scheme}{局部}\Term{splitting}{分裂}，所以 $F$ \Term{local-isomorphism}{局部同构}于 $F'\times F''$，从而\Term{locally-constant}{局部常值}。一般地，取满\Term{homomorphism}{同态} $\varphi:L''\to M''$，其中 $L''$ 为\Term{finite-free-module}{有限型自由模}；$\varphi$ 的核 $R$ 为\Term{finite-type}{有限型}。置
\[
L=L''_T\mathbin{\times}_{M''_T}F=L''_T\mathbin{\times}_{F''}F.
\]
一方面，$L$ 是 $L''_T$ 被 $G$ 的\Term{group-extension}{扩张}，故由上述结论\Term{locally-constant}{局部常值}；另一方面，有\Term{exact-sequence}{正合列} $0\to R_T\to L\to F\to0$，由 (ii) 得 $F$ \Term{locally-constant}{局部常值}。
\end{proof}
从现在起，只考虑\Term{scheme}{概形}的\Term{etale}{平展}\Term{topos}{拓扑斯}。
\begin{SGAresult}{引理}{2.2}
设 $X$ 为一个\Term{scheme}{概形}，$F$ 为 $X$ 上的一个\Term{locally-constant}{局部常值}\Term{sheaf}{层}。则 $F$ 由某个\Term{etale}{平展}的 $Y/X$ 表示。若此外 $F$ 的\Term{fiber}{纤维}\Term{finite}{有限}（分别为\Term{finite}{有限}且\Term{nonempty}{非空}），则 $Y$ 是一个\Term{etale-covering}{平展覆盖}（分别为\Term{surjectivity}{满射}的\Term{etale-covering}{平展覆盖}）。
\end{SGAresult}
\begin{proof}
\Term{descent}{下降}（$F$ 的\Term{fiber}{纤维}\Term{finite}{有限}时见 SGA 1 第 IX 讲 4.1；一般情形见 SGA 3 第 X 讲 5.4）。
\end{proof}
\begin{SGAresult}{定义}{2.3}
若对每个\Term{affine-open}{仿射开集} $U\subset X$，存在一个将 $U$ 写为\Term{finite}{有限}个\Term{locally-closed}{局部闭}、\Term{reduced}{既约}的\Term{constructible}{可构造}子\Term{scheme}{概形} $U_i$（EGA $0_{\mathrm{III}}$ 9.1.2）之并的\Term{resolution}{分解}，使 $F$ 在每个 $U_i$ 上诱导的\Term{sheaf}{层}为\Term{locally-constant}{局部常值}，\Term{local-scheme}{局部}取值\Term{finite}{有限}（分别为\Term{finite}{有限}、\Term{finite-presentation}{有限表现}），就称\Term{set}{集合}\Term{sheaf}{层}（分别为群\Term{sheaf}{层}、$A$-模\Term{sheaf}{层}）为\Term{constructible}{可构造}的。\footnote{不假设 $A$ \Term{noetherian}{诺特}时，为满足\Term{homology}{同调}\Term{algebra}{代数}的需要，最好在此处不只要求 $A$-模 $M$ 为\Term{finite-presentation}{有限表现}，而要求它有一个由\Term{finite-type}{有限型}自由 $A$-模组成的\Term{left-resolution}{左消解}。这时 (2.3) 中引入的概念应当称为：$F$ 是 $1$-\Term{constructible}{可构造}的（更一般地，对每个整数 $n\ge0$，都可以自然地定义 $n$-\Term{constructible}{可构造}性）。在本讲中，我们暂且保留 (2.3) 的术语\Term{form}{形式}；它主要在 $A$ \Term{noetherian}{诺特}时是合理的，此时与刚才指出的概念一致。非\Term{noetherian}{诺特}情形中的\Term{finiteness}{有限性}概念，参看 SGA 6 第 I 讲。}
\end{SGAresult}
\textbf{2.3.1.} 对\Term{local-scheme}{局部}取值的\Term{finiteness}{有限性}假设\Term{equivalence}{等价}于：对每个\Term{geometric-point}{几何点} $\xi$，\Term{fiber}{纤维} $F_\xi$ \Term{finite}{有限}（分别为\Term{finite}{有限}、\Term{finite-presentation}{有限表现}）。由 2.1 (i) 立即得到，群\Term{sheaf}{层} $F$ \Term{constructible}{可构造}，当且仅当其\Term{underlying-set}{底集}合\Term{sheaf}{层}\Term{constructible}{可构造}。

须注意：若 $F$ 是\Term{abelian-sheaf}{阿贝尔层}，则它作为群\Term{sheaf}{层}（或\Term{set}{集合}\Term{sheaf}{层}）\Term{constructible}{可构造}，当且仅当它作为 $\mathbb Z$-模\Term{sheaf}{层}\Term{constructible}{可构造}，并且\Term{fiber}{纤维}\Term{finite}{有限}。例如，$\mathbb Z_X$ 作为 $\mathbb Z$-模\Term{constructible}{可构造}，但作为群\Term{sheaf}{层}则并非如此。
\SourcePage{14}
\begin{SGAresult}{命题}{2.4}
(i) 设 $X$ \Term{quasi-compact}{拟紧}且\Term{quasi-separated}{拟分离}。则\Term{sheaf}{层}（分别为群\Term{sheaf}{层}、$A$-模\Term{sheaf}{层}）$F$ \Term{constructible}{可构造}，当且仅当存在 $X$ 到\Term{locally-closed}{局部闭}\Term{constructible}{可构造}部分之并的\Term{resolution}{分解}，使 $F$ 在每个 $X_i$ 上成为\Term{locally-constant}{局部常值}且\Term{finite}{有限}（分别为\Term{finite}{有限}、\Term{finite-presentation}{有限表现}）的\Term{sheaf}{层}。

(ii) 为验证 $F$ \Term{constructible}{可构造}，只须对 $X$ 的某个给定\Term{open-cover}{开覆盖}中的 $U_i$，验证 $F|U_i$ \Term{constructible}{可构造}。

(iii) 设 $f:Y\to X$ 是\Term{morphism}{态射}。若 $F$ \Term{constructible}{可构造}，则 $f^*F$ 也\Term{constructible}{可构造}。

(iv) \Term{constructible}{可构造}\Term{sheaf}{层}的\Term{fiber}{纤维}\Term{nonempty}{非空}（分别为非零）的点构成一个\Term{local-scheme}{局部}\Term{constructible}{可构造}\Term{subset}{子集}（EGA $0_{\mathrm{III}}$ 9.1.2）。

(v) 设 $X$ \Term{locally-noetherian}{局部诺特}。则 $X$ 上的\Term{set}{集合}\Term{sheaf}{层}（分别为……）$F$ \Term{constructible}{可构造}，当且仅当对每个 $x\in X$，在 $x$ 的\Term{closure}{闭包} $\overline{\{x\}}$ 中存在一个\Term{nonempty}{非空}\Term{open-subset}{开集} $U$，使 $F$ 在 $U$ 上诱导一个取\Term{finite}{有限}值（分别为……）的\Term{locally-constant}{局部常值}\Term{sheaf}{层}。
\end{SGAresult}
\begin{proof}
显然，若 $f:Y\to X$ 是\Term{morphism}{态射}，且存在将 $X$ 写为\Term{locally-closed}{局部闭}\Term{constructible}{可构造}部分 $X_i$ 之并的\Term{resolution}{分解}，使 $F$ 在每个 $X_i$ 上\Term{locally-constant}{局部常值}，则对 $Y$ 与 $f^*F$ 也有这样的\Term{resolution}{分解}（参看 EGA IV 1.8.2）。这证明了 (i) 中的 $\Leftarrow$。现设 $X$ \Term{quasi-compact}{拟紧}且\Term{quasi-separated}{拟分离}，$\{U_1,\ldots,U_n\}$ 为 $X$ 的一个\Term{affine-open-cover}{仿射开覆盖}，且每个 $F|U_i$ 都\Term{constructible}{可构造}。为得到 (i) 中的\Term{resolution}{分解}，只须对 $U_n$ 和 $Y=X-U_n$（赋予诱导的\Term{reduced}{既约}结构）分别得到这种\Term{resolution}{分解}；由“$X$ \Term{quasi-separated}{拟分离}”这一假设，$U_n$ 和 $Y$ 在 $X$ 中都\Term{constructible}{可构造}。由\Term{constructible}{可构造}性的定义，$U_n$ 存在这种\Term{resolution}{分解}；而 $Y$ 是 $n-1$ 个\Term{affine-open}{仿射开集} $V_i=Y\cap U_i$（$i=1,\ldots,n-1$）之并，且 $F|V_i$ \Term{constructible}{可构造}，故对 $n$ 归纳即得 (i)。同样的论证证明 (ii)：事实上，假设意味着对每个“充分小”的\Term{affine-open}{仿射开集} $U$，$F|U$ \Term{constructible}{可构造}；而任意\Term{affine-open}{仿射开集} $V$ 都可由\Term{finite}{有限}个这样的\Term{open-subset}{开集}\Term{covering}{覆盖}。(ii) 表明\Term{constructible}{可构造}性在 $X$ 上是\Term{local-nature}{局部性质}，(iii)、(iv) 随即得到。(v) 中的 $\Rightarrow$ 是直接的，而且不依赖\Term{noetherian}{诺特}假设。若 $X$ \Term{noetherian}{诺特}，则用通常的“\Term{noetherian-induction}{诺特归纳}法”证明 $\Leftarrow$。
\end{proof}
\begin{SGAresult}{命题}{2.5}
设 $X$ \Term{quasi-compact}{拟紧}且\Term{quasi-separated}{拟分离}，$F$ 是 $X$ 上的\Term{constructible}{可构造}群\Term{sheaf}{层}或 $A$-模\Term{sheaf}{层}。则 $F$ 存在一个\Term{finite}{有限}\Term{filtered}{滤过}，其逐次商形如 $i_!G$，其中 $i:U\to X$ 是一个\Term{locally-closed}{局部闭}\Term{constructible}{可构造}部分的包含，$G$ 在 $U$ 上\Term{locally-constant}{局部常值}且\Term{constructible}{可构造}。若 $X$ \Term{noetherian}{诺特}，则可使这些 $U$ 均\Term{irreducible}{不可约}。
\end{SGAresult}
\begin{proof}
由 2.4 (i)，存在 $X$ 到\Term{finite}{有限}个\Term{locally-closed}{局部闭}\Term{constructible}{可构造}部分 $X_i$ 之并的\Term{resolution}{分解}，使每个 $F|X_i$ 都\Term{locally-constant}{局部常值}且\Term{finite}{有限}（分别为\Term{finite-presentation}{有限表现}）。写 $X_i=U_i\cap\complement V_i$，其中 $U_i,V_i$ 均为\Term{constructible}{可构造}\Term{open-subset}{开集}。
\SourcePage{15}
设 $N$ 为 $\{U_i,V_i\}$ 所生成的子\Term{topology}{拓扑} $\mathcal T$ 中 $X$ 的\Term{open-subset}{开集}的数目。对 $N$ 归纳：若 $W$ 为 $\mathcal T$ 的一个极小\Term{nonempty}{非空}元素，显然 $F|W$ \Term{locally-constant}{局部常值}且为\Term{finite-presentation}{有限表现}。设 $i:W\to X$ 为\Term{inclusion-morphism}{包含态射}，$Y=X-W$。有\Term{exact-sequence}{正合列}
\[
0\longrightarrow i_!i^*F\longrightarrow F\longrightarrow F|Y\longrightarrow0,
\]
于是归结到 $Y$ 与\Term{sheaf}{层} $F|Y$ 的同一断言，此时可应用归纳假设。
\end{proof}
\begin{SGAresult}{命题}{2.6}
设 $A$ 为\Term{noetherian-ring}{诺特环}。

(i) \Term{constructible}{可构造}的 $A$-模\Term{sheaf}{层}或\Term{set}{集合}\Term{sheaf}{层}的\Term{finite}{有限}\Term{projective-limit}{投射极限}（分别为\Term{finite}{有限}\Term{inductive-limit}{归纳极限}）仍\Term{constructible}{可构造}。特别地，\Term{constructible}{可构造}的 $A$-模\Term{sheaf}{层}之间的\Term{morphism}{态射}的核、\Term{cokernel}{余核}及像均\Term{constructible}{可构造}。

(i bis) \Term{constructible}{可构造}\Term{set}{集合}\Term{sheaf}{层}（分别为群\Term{sheaf}{层}）的\Term{finite}{有限}\Term{projective-limit}{投射极限}\Term{constructible}{可构造}；若 $f:F\to G$ 是\Term{constructible}{可构造}群\Term{sheaf}{层}之间的\Term{morphism}{态射}，则其核、像以及\Term{cokernel}{余核}（若 $\operatorname{Im}(f)$ 为\Term{normal-subgroup}{正规子群}）均\Term{constructible}{可构造}。\Term{constructible}{可构造}\Term{set}{集合}\Term{sheaf}{层}的\Term{finite}{有限}\Term{inductive-limit}{归纳极限}\Term{constructible}{可构造}。

(ii) 设 $0\to M'\to M\to M''\to0$ 为 $A$-模\Term{sheaf}{层}（分别为群\Term{sheaf}{层}）的\Term{exact-sequence}{正合列}，且 $M',M''$ \Term{constructible}{可构造}。则 $M$ \Term{constructible}{可构造}。

(iii) 设 $M_1\to M_2\to M_3\to M_4\to M_5$ 为 $A$-模\Term{sheaf}{层}的\Term{exact-sequence}{正合列}。若 $i=1,2,4,5$ 时 $M_i$ \Term{constructible}{可构造}，则 $i=3$ 时也如此。
\end{SGAresult}
\begin{proof}
(i) \Term{set}{集合}与群的断言留给读者。在 $A$-模\Term{sheaf}{层}的情形，只须证明\Term{constructible}{可构造}的 $A$-模\Term{sheaf}{层}之间的\Term{morphism}{态射} $f:F\to G$ 的核与\Term{cokernel}{余核}\Term{constructible}{可构造}。断言在 $X$ 上是\Term{local-scheme}{局部}的，因而可设 $X$ \Term{affine}{仿射}，从而\Term{quasi-compact}{拟紧}且\Term{quasi-separated}{拟分离}。于是由 2.4 (i) 可归结到 $F,G$ \Term{locally-constant}{局部常值}的情形，再应用 2.1 (ii) 和 $A$ 的\Term{noetherian}{诺特}假设即得。利用 2.1 (iii)，同样证明 (ii)。(iii) 由 (i)、(ii) 立即得到。
\end{proof}
\begin{SGAresult}{命题}{2.7}
设 $X$ 为\Term{quasi-compact}{拟紧}\Term{quasi-separated}{拟分离}\Term{scheme}{概形}，$F$ 为\Term{set}{集合}\Term{sheaf}{层}（分别为 $A$-模\Term{sheaf}{层}）。$F$ \Term{constructible}{可构造}的充要条件是：它\Term{isomorphism}{同构}于一对\Term{morphism}{态射} $H\rightrightarrows G$ 的\Term{cokernel}{余核}，其中 $H,G$ 是由 $X$ 上\Term{finite-presentation}{有限表现}的\Term{etale}{平展}\Term{scheme}{概形}所表示的\Term{set}{集合}\Term{sheaf}{层}（分别为：$F$ \Term{isomorphism}{同构}于某个\Term{homomorphism}{同态} $A_{V,X}\to A_{U,X}$ 的\Term{cokernel}{余核}，其中 $U,V$ 是 $X$ 上\Term{finite-presentation}{有限表现}的\Term{etale}{平展}\Term{scheme}{概形}）。
\end{SGAresult}
\begin{proof}
充分性容易由 2.6 (i bis) 得到。
\SourcePage{16}
设 $F$ 为\Term{constructible}{可构造}\Term{set}{集合}\Term{sheaf}{层}。利用 $X$ 的\Term{etale}{平展}\Term{site}{景}中任意和的存在性，可找到一个\Term{epimorphism}{满态射} $G\to F$，其中 $G$ 由 $X$ 上的一个\Term{etale}{平展}\Term{scheme}{概形}表示；还可设它是\Term{affine}{仿射}\Term{scheme}{概形} $G_i$（$i\in I$）的和，因而在 $S$ 上\Term{separated}{分离}。对指标集 $I$ 的每个\Term{finite}{有限}\Term{subset}{子集} $J$，令 $G_J$ 为 $i\in J$ 时各 $G_i$ 的和，$F_J$ 为它在 $F$ 中的像。由于 $G_J,F$ \Term{constructible}{可构造}，$F_J$ 也\Term{constructible}{可构造}（2.6 (i bis)）；因此，$F$ 与 $F_J$ 在 $x$ 上方的\Term{geometric-point}{几何点}处\Term{fiber}{纤维}相同的那些 $x\in X$ 所成的\Term{set}{集合} $X_J$ 是\Term{local-scheme}{局部}\Term{constructible}{可构造}的（2.4 (iv)）。族 $X_J$ 递增且并为 $X$，而 $X$ \Term{quasi-compact}{拟紧}，故某个 $X_J$ 等于 $X$（EGA IV 1.9.9）。这说明，将 $G$ 换成 $G_J$ 后，可设 $G$ \Term{affine}{仿射}，故在 $X$ 上\Term{separated}{分离}；又因 $X$ \Term{quasi-separated}{拟分离}，它在 $X$ 上\Term{quasi-compact}{拟紧}（EGA IV 1.2.4）。令 $H=G\times_FG$。它是\Term{representable-functor}{可表示}\Term{sheaf}{层} $G\times_SG$ 的子\Term{sheaf}{层}，因此自身也\Term{representable-functor}{可表示}（第 VIII 讲 6.1）。又因 $G$ \Term{constructible}{可构造}（2.6 (i bis)），上述论证仍表明它\Term{quasi-compact}{拟紧}，从而（因它在 $X$ 上\Term{separated}{分离}）在 $X$ 上为\Term{finite-presentation}{有限表现}。这证明了 2.7 的第一条断言；第二条用同样的方法证明（注意不必假设 $A$ \Term{noetherian}{诺特}）。顺便指出，此证明还得到下列推论。
\end{proof}
\begin{SGAresult}{推论}{2.7.1}
设 $X$ 为\Term{scheme}{概形}，$F$ 为 $X$ 上的\Term{etale}{平展}\Term{scheme}{概形}。$X$ 上相应的\Term{etale}{平展}\Term{sheaf}{层}\Term{constructible}{可构造}的充要条件是 $F$ 在 $X$ 上为\Term{finite-presentation}{有限表现}。
\end{SGAresult}
在这里，命题 2.7 的主要用途，是推出\Term{constructible}{可构造}\Term{sheaf}{层}的某些取\Term{limit}{极限}性质。
\begin{SGAresult}{推论}{2.7.2}
设 $X$ 为\Term{quasi-compact}{拟紧}\Term{quasi-separated}{拟分离}\Term{scheme}{概形}。则 $X$ 上的每个\Term{set}{集合}\Term{sheaf}{层}（分别为\Term{ind-p-group}{ind-$\mathfrak p$-群}\Term{sheaf}{层}、$A$-模\Term{sheaf}{层}）都是一个由\Term{constructible}{可构造}\Term{set}{集合}\Term{sheaf}{层}（分别为……）组成的\Term{filtered-inductive-system}{滤过归纳系统}的\Term{inductive-limit}{归纳极限}。
\end{SGAresult}
\begin{proof}
若 $F$ 是\Term{set}{集合}\Term{sheaf}{层}，沿用 2.7 证明中的\Term{signature}{符号}；令 $L$ 为所有二元组 $(J,H')$ 的\Term{set}{集合}，其中 $J$ 是 $I$ 的\Term{finite}{有限}\Term{subset}{子集}，$H'$ 是 $H_J=G_J\times_FG_J$ 的一个\Term{quasi-compact}{拟紧}\Term{open-subset}{开子集}。在 $L$ 上赋予显然的序。由于 $F$ 是各 $F_J$ 的\Term{limit}{极限}，且每个 $F_J$ 是 $\operatorname{Coker}(H'\rightrightarrows G_J)$ 的\Term{limit}{极限}，其中 $H'$ 遍历 $H_J$ 的\Term{quasi-compact}{拟紧}\Term{open-subset}{开子集}，因此 $F$ 是\Term{inductive-system}{归纳系统} $F_\ell=\operatorname{Coker}(H'\rightrightarrows F_J)$（$\ell=(H',J)\in L$）的\Term{inductive-limit}{归纳极限}，\Term{set}{集合}情形即得。$A$-模\Term{sheaf}{层}的情形本质上相同：先将 $F$ 表成\Term{homomorphism}{同态} $A_{V,X}\to A_{U,X}$ 的\Term{cokernel}{余核}，其中 $U,V$ 是 $X$ 上\Term{separated}{分离}的\Term{etale}{平展}\Term{scheme}{概形}。\Term{ind-p-group}{ind-$\mathfrak p$-群}的情形需要稍加注意。沿用上述\Term{signature}{符号}，令 $L(G_J)$ 为\Term{set}{集合}\Term{sheaf}{层} $G_J$ 所生成的\Term{free-group}{自由群}\Term{sheaf}{层}，$N_J$ 为从 $G_J\to F$ 得到的\Term{homomorphism}{同态} $L(G_J)\to F$ 的核；令 $L$ 为所有二元组 $(J,N')$ 的\Term{set}{集合}，
\SourcePage{17}
其中 $N'$ 是 $N_J$ 的一个“\Term{quasi-compact}{拟紧}”\Term{set}{集合}子\Term{sheaf}{层}，即存在\Term{epimorphism}{满态射} $P\to N'$，且 $P$ 由一个\Term{quasi-compact}{拟紧} $X$-\Term{scheme}{概形}表示。定义 $(J,N')\le(J_1,N'_1)$ 当且仅当 $J\subset J_1$，且\Term{canonical-homomorphism}{典范同态} $L(G_J)\to L(G_{J_1})$ 将 $N'$ 映入 $N'_1$。显然 $L$ 是\Term{filtered-category}{滤过的}，并且令 $\ell=(J,N')$ 时，$F_\ell$ 为 $L(G_J)$ 模掉 $N'$ 所生成的不变群子\Term{sheaf}{层}的\Term{quotient-group}{商群}\Term{sheaf}{层}，就得到一个以 $L$ 为指标的群\Term{sheaf}{层}\Term{inductive-system}{归纳系统}。同样显然，$F$ \Term{isomorphism}{同构}于各 $F_\ell$ 的\Term{inductive-limit}{归纳极限}。还须证明存在 $L$ 的一个\Term{cofinal}{余尾}\Term{subset}{子集} $L'$，使 $\ell\in L'$ 时，$F_\ell$ 是一个\Term{constructible}{可构造}的 $\mathfrak p$-群\Term{sheaf}{层}。为此，只须证明：对每个 $J$，给定\Term{quasi-compact}{拟紧}\Term{set}{集合}子\Term{sheaf}{层} $N'_0$，可找到 $N_J$ 的一个\Term{quasi-compact}{拟紧}\Term{set}{集合}子\Term{sheaf}{层} $N'$ 包含 $N'_0$，使对应的 $F_\ell$ 是\Term{constructible}{可构造}的 $\mathfrak p$-群\Term{sheaf}{层}。

先注意，只须使 $F_\ell$ 的\Term{fiber}{纤维}都是\Term{finite}{有限} $\mathfrak p$-群；这时 $F_\ell$ 自动\Term{constructible}{可构造}，读者容易验证。另一方面，读者同样可验证：$F_\ell$ 在 $x$ 上方某\Term{geometric-point}{几何点}处的\Term{fiber}{纤维}为\Term{finite}{有限} $\mathfrak p$-群的那些点 $x\in X$ 所成的\Term{set}{集合} $X_{N'}$，是\Term{local-scheme}{局部}\Term{constructible}{可构造}的。而且 $X$ 的这些部分 $X_{N'}$ 随 $N'$ 递增；最后，它们的并为整个 $X$。最后这一点由下述事实得到：固定 $J$，$\ell=(N',J)$ 时各 $F_\ell$ 的\Term{inductive-limit}{归纳极限}，是\Term{set}{集合}子\Term{sheaf}{层} $F_J$ 所生成的 $F$ 的群子\Term{sheaf}{层} $\overline F_J$；$F_J$ \Term{quasi-compact}{拟紧}，所以 $\overline F_J$ 的\Term{fiber}{纤维}都是\Term{finite}{有限} $\mathfrak p$-群。另外，在一个\Term{finitely-generated}{有限生成}的\Term{free-group}{自由群}中，每个\Term{finite-index}{有限指数}\Term{subgroup}{子群}都是\Term{finitely-generated}{有限生成}\Term{free-group}{自由群}；故若将该\Term{subgroup}{子群}写成其\Term{finite}{有限}\Term{subset}{子集}的\Term{filtered}{滤过}并，则其中一个\Term{finite}{有限}\Term{subset}{子集}就生成整个\Term{subgroup}{子群}。由此容易得到，$X_{N'}$ 的并确为 $X$。再次应用 EGA IV 1.9.9，得到某个 $X_{N'}=X$，完成 2.7.2 的证明。
\end{proof}
\textbf{2.7.2.1.} 群\Term{scheme}{概形}情形的这个有些繁琐的证明，显示出我们缺少一套方便的基本性质，来处理\Term{fiber}{纤维}未必\Term{finite}{有限}的群\Term{sheaf}{层}的\Term{constructible}{可构造}性条件。我们为此向读者致歉，并请读者在需要时自行填补这一空缺。聊以慰藉的是，在 $X$ \Term{noetherian}{诺特}时，2.9 (iii) 将给出一个更简单的证明。
\begin{SGAresult}{推论}{2.7.3}
设 $X$ 为\Term{quasi-compact}{拟紧}\Term{quasi-separated}{拟分离}\Term{scheme}{概形}，$F$ 为\Term{constructible}{可构造}\Term{set}{集合}\Term{sheaf}{层}（分别为\Term{ind-finite-group}{ind-有限群}\Term{sheaf}{层}、$A$-模\Term{sheaf}{层}）。则\Term{functor}{函子} $\operatorname{Hom}(F,G)$ 与变量\Term{set}{集合}\Term{sheaf}{层}（分别为……）$G$ 的\Term{inductive-limit}{归纳极限}交换。在 $F$ 为 $A$-模\Term{sheaf}{层}且 $A$ \Term{noetherian}{诺特}的情形，关于 $A$-模 $G$ 的\Term{functor}{函子} $\operatorname{Ext}^i(X;F,G)$ 也与\Term{inductive-limit}{归纳极限}交换。
\end{SGAresult}
\SourcePage{18}
这由 2.7 和第 VII 讲 3.3 通过通常的论证得到，留给读者。同样，将 2.7 与第 VII 讲 5 的结果结合，容易得到：
\begin{SGAresult}{推论}{2.7.4}
设 $I$ 为\Term{filtered-category}{滤过范畴}，$i\mapsto X_i:I^\circ\to(\mathrm{Sch})$ 是将\Term{arrow}{箭头}变为\Term{affine-morphism}{仿射态射}、\Term{object}{对象}变为\Term{quasi-compact}{拟紧}\Term{quasi-separated}{拟分离}\Term{scheme}{概形}的\Term{functor}{函子}。对每个\Term{scheme}{概形} $Y$，令 $F_c(Y)$ 为 $Y$ 上\Term{constructible}{可构造}\Term{set}{集合}\Term{sheaf}{层}（分别为\Term{ind-p-group}{ind-$\mathfrak p$-群}\Term{sheaf}{层}、$A$-模\Term{sheaf}{层}）的\Term{category}{范畴}。则有\Term{category-equivalence}{范畴等价}
\[
F_c(X)\xleftarrow{\approx}\varinjlim_iF_c(X_i),
\]
其中 $X=\varprojlim_iX_i$（第 VII 讲 5），右端表示\Term{fibred-category}{纤维范畴}的\Term{limit}{极限}\Term{category}{范畴} $\varinjlim$，其定义见\SGAsourcegap{原重排本引用空缺，补源原版第 IX 讲第 15 页此处也留空}。特别地，$X$ 上每个\Term{constructible}{可构造}\Term{set}{集合}\Term{sheaf}{层}（分别为……）都\Term{isomorphism}{同构}于某个 $X_i$ 上同类\Term{sheaf}{层}的\Term{inverse-image}{逆像}。
\end{SGAresult}
\begin{SGAresult}{命题}{2.8}
设 $f:X\to Y$ 为\Term{surjectivity}{满射}且\Term{local-scheme}{局部}\Term{finite-presentation}{有限表现}的\Term{morphism}{态射}，$F$ 为 $Y$ 上的\Term{set}{集合}\Term{sheaf}{层}（分别为……）。则 $F$ \Term{constructible}{可构造}，当且仅当 $f^*(F)$ \Term{constructible}{可构造}。
\end{SGAresult}
我们将使用下述引理。
\begin{SGAresult}{引理}{2.8.1}
设 $f:X\to Y$ 为\Term{surjectivity}{满射}且\Term{local-scheme}{局部}\Term{finite-presentation}{有限表现}的\Term{morphism}{态射}，$Y$ \Term{quasi-compact}{拟紧}且\Term{quasi-separated}{拟分离}。则存在 $Y$ 到在 $Y$ 上为\Term{finite-presentation}{有限表现}的子\Term{scheme}{概形} $Y_i$ 的一个\Term{finite}{有限}分划（特别地，每个 $Y_i$ 在 $X$ 中\Term{constructible}{可构造}）；对每个 $i$，有\Term{finite}{有限}\Term{surjectivity}{满射}
\[
Y_i''\xrightarrow{g_i}Y_i'\xrightarrow{h_i}Y_i,
\]
其中 $h_i$ \Term{etale}{平展}，$g_i$ \Term{locally-free}{局部自由}（即除了\Term{finite}{有限}这一条件外，还\Term{flat}{平坦}且为\Term{finite-presentation}{有限表现}）且\Term{radicial}{根性}；最后还有一个 $Y$-\Term{morphism}{态射} $Y_i''\to X$。\footnote{这一命题也见于 EGA IV 17.16.4。}
\end{SGAresult}
\begin{proof}
利用 $Y$ \Term{quasi-compact}{拟紧}且\Term{quasi-separated}{拟分离}，容易先归结到 $Y$ \Term{quasi-affine-morphism}{拟仿射}因而\Term{separated}{分离}的情形（取 $Y$ 的\Term{finite}{有限}\Term{affine-open-cover}{仿射开覆盖} $(U_i)_{1\le i\le n}$，考虑 $Y_i=U_i-(U_i\cap\bigcup_{j<i}U_j)$），再用同样的论证归结到 $Y$ \Term{affine}{仿射}。\Term{scheme}{概形} $X$ 是\Term{affine-open}{仿射开集} $X_i$ 的并；$X_i$ 在 $Y$ 中的像 $T_i$ \Term{constructible}{可构造}（EGA IV 1.8.4），由于 $f$ \Term{surjectivity}{满射}，$T_i$ 的并为 $X$，所以 $X$ 是\Term{finite}{有限}个 $T_i$ 的并（EGA IV 1.9.9）。于是，将 $X$ 换成相应 $X_i$ 的和，可设 $X$ \Term{affine}{仿射}。再由标准的取\Term{limit}{极限}方法（EGA IV 8.1.2 c)），可进一步设 $Y$ \Term{noetherian}{诺特}。通常的\Term{noetherian-induction}{诺特归纳}法和取\Term{limit}{极限}方法（EGA IV 8.1.2 a)）将我们归结到 $Y$ 是一个域 $k$ 的谱的情形。这时 $X\ne\varnothing$，故 $X$ 有一个\Term{closed-point}{闭点} $x$，对应 $k$ 的一个\Term{finite-extension}{有限扩张} $k'$，它又是 $k$ 的某个\Term{separable-extension}{可分扩张} $k'_s$ 的\Term{radicial}{根性}\Term{group-extension}{扩张}。
\SourcePage{19}
取 $Y'=\operatorname{Spec}(k'_s)$、$Y''=\operatorname{Spec}(k')$，即完成 2.8.1 的证明。
\end{proof}
\begin{proof}[命题 2.8 的证明]
必要性是直接的（2.4 (iii)），现在证明充分性。由上述引理，可归结到 $f$ 形如 $gh$，其中 $g,h$ 满足引理对 $g_i,h_i$ 所述的条件。于是\Term{separable}{可分}别设 $f$ 为\Term{finite}{有限}\Term{radicial}{根性}\Term{morphism}{态射}（由第 VIII 讲 1.1，这是平凡的情形），或 $f$ 为\Term{finite}{有限}\Term{etale}{平展}\Term{morphism}{态射}。还要注意，利用 $f^*(F)$ \Term{constructible}{可构造}这一假设和\Term{constructible}{可构造}性的定义，将 $X$ 换成一个在 $X$ 上为\Term{finite-presentation}{有限表现}且\Term{structural-morphism}{结构态射}为\Term{surjectivity}{满射}的 $X$-\Term{prescheme}{预概形} $X'=\coprod_iX_i$ 后，可设 $f^*(F)$ 已经\Term{locally-constant}{局部常值}。而当 $f$ \Term{etale}{平展}且\Term{surjectivity}{满射}时，这意味着 $F$ \Term{locally-constant}{局部常值}。由于 $f$ \Term{surjectivity}{满射}，对 $f^*(F)$ 的\Term{fiber}{纤维}所作的假设对 $F$ 的\Term{fiber}{纤维}仍成立，所以 $F$ 确为\Term{constructible}{可构造}，2.8 得证。
\end{proof}
\begin{SGAresult}{命题}{2.9}
设 $X$ 为\Term{noetherian}{诺特}\Term{scheme}{概形}，$A$ 为\Term{noetherian-ring}{诺特环}，$\mathfrak p$ 为一个素数\Term{set}{集合}。

(i) $X$ 上\Term{set}{集合}\Term{sheaf}{层}（分别为\Term{ind-p-group}{ind-$\mathfrak p$-群}\Term{sheaf}{层}、$A$-模\Term{sheaf}{层}）的\Term{category}{范畴}\Term{locally-noetherian}{局部诺特}，即有一个由\Term{noetherian}{诺特}\Term{object}{对象}组成的\Term{generator}{生成元}集。

(ii) $X$ 上的\Term{set}{集合}\Term{sheaf}{层}（分别为\Term{ind-p-group}{ind-$\mathfrak p$-群}\Term{sheaf}{层}、$A$-模\Term{sheaf}{层}）$F$ \Term{constructible}{可构造}，当且仅当它是\Term{noetherian}{诺特}\Term{object}{对象}。若 $F$ 是 $A$-模\Term{sheaf}{层}，则 $F$ \Term{constructible}{可构造}，当且仅当它是\Term{finite}{有限}个形如 $A_{U/X}$ 的\Term{sheaf}{层}的\Term{direct-sum}{直和}的商，其中 $U\to X$ \Term{etale}{平展}且为\Term{finite-type}{有限型}（第 IV 讲 2.4 的\Term{signature}{符号}：$A_{U/X}=A_U$）。

(iii) \Term{set}{集合}\Term{sheaf}{层}（分别为……）$F$ 的\Term{constructible}{可构造}子\Term{sheaf}{层}构成一个\Term{inductive-system}{归纳系统}，且 $F$ 是这些子\Term{sheaf}{层}的\Term{inductive-limit}{归纳极限}。
\end{SGAresult}
\begin{proof}
\Term{set}{集合}情形的断言留给读者。先处理 $A$-模\Term{sheaf}{层}。为证明 (i)，须找到一个\Term{generator}{生成元}集，其元素都是\Term{noetherian}{诺特}\Term{object}{对象}。而形如 $A_{U/X}$（$U\to X$ \Term{etale}{平展}且为\Term{finite-type}{有限型}）的\Term{sheaf}{层}构成一个\Term{generator}{生成元}集，且容易看出它们\Term{constructible}{可构造}，所以只须证明下述引理。
\end{proof}
\begin{SGAresult}{引理}{2.10}
设 $X$ \Term{noetherian}{诺特}。则每个\Term{constructible}{可构造}的 $A$-模都是\Term{noetherian}{诺特}\Term{object}{对象}。
\end{SGAresult}
\begin{proof}[引理的证明]
由\Term{noetherian-induction}{诺特归纳}法，可设对 $X$ 的每个不同于 $X$ 的子\Term{scheme}{概形} $X'$，引理都成立。取 $X$ 的一个\Term{nonempty}{非空}\Term{open-subset}{开集} $X_0$，使 $F$ 在 $X_0$ 上诱导一个\Term{locally-constant}{局部常值}\Term{sheaf}{层} $F_0$。

现设 $\{F_i\}$，$i\in\mathbb N$，为 $F_0$ 的一列递增子\Term{sheaf}{层}。对 $X_0$ 某个\Term{irreducible-component}{不可约分支}的泛\Term{geometric-point}{几何点} $P$，由于 $F_0$ 在 $P$ 处的\Term{fiber}{纤维}为\Term{finite-type}{有限型}，\Term{fiber}{纤维}序列 $(F_i)_P$ \Term{stable}{稳定}，可设它为常值序列。
\SourcePage{20}
设 $Q$ 为 $X_0$ 的一个\Term{geometric-point}{几何点}，是 $P$ 的\Term{specialization}{特殊化}（第 VIII 讲 7.2）。因 $F_0$ 在 $X_0$ 上\Term{locally-constant}{局部常值}，立即可见\Term{specialization}{特殊化}\Term{morphism}{态射} $F_{0Q}\to F_{0P}$（第 VIII 讲 7）是\Term{bijection}{双射}，故对每个 $i$，$(F_i)_Q\to(F_i)_P$ 都是\Term{injectivity}{单射}。

设 $V/X$ \Term{etale}{平展}且带 $P$ 标点，$s_1,\ldots,s_n\in F_1(V)$ 是生成 $(F_1)_P$ 的\Term{section}{截面}。则对每个 $i$，以及在 $V$ 上方作为 $P$ 的\Term{specialization}{特殊化}的每个\Term{geometric-point}{几何点} $Q$，$s_i$ 都生成 $(F_i)_Q$。因此\Term{sheaf}{层}序列 $F_i$ 在 $P$ 的像的某个\Term{neighborhood}{邻域}内为常值序列，记此\Term{neighborhood}{邻域}为 $X'\ne\varnothing$。令 $Y=X-X'$，于是只须证明 $F_i|Y$ 的序列\Term{stable}{稳定}；这由\Term{noetherian-induction}{诺特归纳}假设得到。
\end{proof}
\begin{proof}[命题 2.9 证明的续]
对于 $A$-模，(ii) 现在是直接的：由于 $A_{U/X}$ 是\Term{noetherian}{诺特}\Term{generator}{生成元}，显然
\[
F\text{ 诺特}\ \Longleftrightarrow\ F\text{ 是有限个 }A_{U/X}\text{ 的直和的商}.
\]
而 $A_{U/X}$ \Term{constructible}{可构造}，故结合 2.6，$F$ \Term{noetherian}{诺特}推出 $F$ \Term{constructible}{可构造}；反向蕴含已由 2.10 证明，因此 (ii) 成立。最后，$A$-模情形的 (iii) 是\Term{locally-noetherian}{局部诺特}\Term{category}{范畴}的性质（参看 [1]，第 IV 章）。

现在处理\Term{ind-p-group}{ind-$\mathfrak p$-群}。先注意，\Term{constructible}{可构造}群\Term{sheaf}{层} $F$ 必为\Term{noetherian}{诺特}\Term{object}{对象}。事实上，由于\Term{fiber}{纤维}\Term{finite}{有限}，可使用与 $A$-模相同的论证。现在直接证明 (iii)，它将蕴含 (i)：这实际上表明\Term{constructible}{可构造}\Term{sheaf}{层}构成一个\Term{generator}{生成元}集。而且，每个\Term{noetherian}{诺特}\Term{sheaf}{层}这时必\Term{constructible}{可构造}，故 (ii) 也成立。

为证明 (iii)，显然只须证明：\Term{ind-p-group}{ind-$\mathfrak p$-群}\Term{sheaf}{层} $F$ 中由\Term{finite}{有限}个\Term{section}{截面} $s_i\in F(U_i)$（$U_i\to X$ \Term{etale}{平展}且为\Term{finite-type}{有限型}）生成的群子\Term{sheaf}{层} $S$ 是\Term{constructible}{可构造}的。\footnote{回忆此\Term{sheaf}{层}的定义：令 $L_{U/X}$ 为在群\Term{category}{范畴}中表示\Term{functor}{函子} $F\mapsto F(U)$ 的\Term{sheaf}{层}。\Term{section}{截面} $s_i\in F(U_i)$ 诱\Term{derived-morphism}{导态射} $L_{U_i/X}\to F$，$S$ 是和\Term{morphism}{态射} $\coprod_iL_{U_i/X}\to F$ 的像。这是“包含”各\Term{section}{截面} $s_i$ 的最小子\Term{sheaf}{层}。}（由于\Term{constructible}{可构造}\Term{sheaf}{层}是\Term{noetherian}{诺特}\Term{object}{对象}，它可由\Term{finite}{有限}个\Term{section}{截面}生成。）由\Term{noetherian-induction}{诺特归纳}法，可设对 $X$ 的每个不同于 $X$ 的闭子\Term{scheme}{概形} $Y$，此断言都成立。由\Term{section}{截面}生成的子\Term{sheaf}{层}这一构造与\Term{sheaf}{层}的\Term{inverse-image}{逆像}交换：事实上，令 $L_{U/X}$ 为在群\Term{category}{范畴}中表示\Term{functor}{函子} $F\mapsto F(U)$ 的\Term{sheaf}{层}。若 $f:Y\to X$，则
\[
f^*(L_{U/X})=L_{Y\times_XU/Y}.
\]
因为 $S$ 是各 $L_{U_i/X}$ 之和的像，且 $f^*$ 与和交换，显然 $f^*S$ 是 $f^*F$ 中由\Term{section}{截面} $f^*(s_i)$ 生成的子\Term{sheaf}{层}。

所以可设 $S$ 在每个不同于 $X$ 的闭集 $Y$ 上诱导一个\Term{constructible}{可构造}\Term{sheaf}{层}，因而只须（2.4 (i)）证明 $S$ 在 $X$ 的某个适当\Term{nonempty}{非空}\Term{open-subset}{开集}上诱导一个\Term{constructible}{可构造}\Term{sheaf}{层}。将 $X$ 换成充分小的\Term{open-subset}{开集}，可设每个 $U_i$ 都是 $X$ 的\Term{etale-covering}{平展覆盖}。
\SourcePage{21}
只须证明此时 $S$ 在某个适当\Term{nonempty}{非空}\Term{open-subset}{开集}上\Term{locally-constant}{局部常值}且\Term{fiber}{纤维}\Term{finite}{有限}。这是 $X$ 上关于\Term{etale}{平展}\Term{topology}{拓扑}的\Term{local-scheme}{局部}断言，所以可设每个 $U_i$ 都\Term{completely-split}{完全分裂}，例如 $U_i=\coprod^{n_i}X$（$X$ 的 $n_i$ 个副本之和）。设 $s_{ij}$（$j=1,\ldots,n_i$）为 $F(X)$ 的\Term{section}{截面}，使 $s_i=(s_{i1},\ldots,s_{in_i})$。显然，$S$ 正是 $F$ 中由 $s_{ij}\in F(X)$ 生成的子\Term{sheaf}{层}。因 $F(X)$ 是\Term{ind-p-group}{ind-$\mathfrak p$-群}（1.5），各 $s_{ij}$ 生成一个\Term{finite}{有限}\Term{subgroup}{子群}（1.3）。所以可设 $s_i$ 的\Term{set}{集合}是 $F(X)$ 的一个\Term{finite}{有限}\Term{subgroup}{子群}；这时 $S$ 就等于这些 $s_i$ 所生成的\Term{set}{集合}\Term{sheaf}{层}。由 2.6 (i)，它作为\Term{set}{集合}\Term{sheaf}{层}\Term{constructible}{可构造}，因而作为群\Term{sheaf}{层}也\Term{constructible}{可构造}。证毕。
\end{proof}
下面两个命题利用\Term{specialization}{特殊化}\Term{homomorphism}{同态}（第 VIII 讲 7.7），给出\Term{sheaf}{层}\Term{locally-constant}{局部常值}或\Term{constructible}{可构造}的判据。
\begin{SGAresult}{命题}{2.11}
设 $X$ 为\Term{scheme}{概形}，$F$ 为\Term{constructible}{可构造}\Term{set}{集合}\Term{sheaf}{层}（分别为\Term{ind-finite-group}{ind-有限群}\Term{sheaf}{层}、$A$-模\Term{sheaf}{层}）；在\Term{set}{集合}情形，另设 $F$ 的\Term{fiber}{纤维}\Term{finite}{有限}。设 $x\in X$，$\bar x$ 为 $x$ 上方的\Term{geometric-point}{几何点}。$F$ 在 $x$ 的某\Term{neighborhood}{邻域}内\Term{locally-constant}{局部常值}的充要条件是：对每个为 $\bar x$ 的\Term{generalization}{一般化}的\Term{geometric-point}{几何点} $\bar x'$，以及每个\Term{specialization}{特殊化}\Term{arrow}{箭头} $\bar x'\to\bar x$（第 VIII 讲 7.2），\Term{specialization}{特殊化}\Term{homomorphism}{同态}（第 VIII 讲 7.7）$F_{\bar x}\to F_{\bar x'}$ 都是\Term{isomorphism}{同构}。
\end{SGAresult}
\begin{proof}
群\Term{sheaf}{层}情形是\Term{set}{集合}\Term{sheaf}{层}情形的特例，所以只处理后者；模\Term{sheaf}{层}的情形本质上相同。条件的必要性是直接的，只须证明充分性。由于\Term{fiber}{纤维} $I=F_{\bar x}$ \Term{finite}{有限}，将 $X$ 换成 $X$ 上某个适当的\Term{etale}{平展}\Term{scheme}{概形} $X'$ 后，可设存在由\Term{constant-sheaf}{常层} $I_X$ 到 $F$ 的\Term{homomorphism}{同态} $u:I_X\to F$，它在 $\bar x$ 处的\Term{fiber}{纤维}上\Term{induction-isomorphism}{诱导同构}。利用这两个\Term{sheaf}{层}的\Term{constructible}{可构造}性及 2.6，容易看出：在 $z$ 上方某\Term{geometric-point}{几何点} $\bar z$ 的\Term{fiber}{纤维}上诱导的\Term{homomorphism}{同态}为\Term{bijection}{双射}的那些点 $z\in X$，构成 $X$ 的一个\Term{local-scheme}{局部}\Term{constructible}{可构造}部分 $Z$。假设立即意味着它包含 $x$ 的所有\Term{generalization}{一般化}，因而它是 $x$ 的\Term{neighborhood}{邻域}（EGA IV 1.10.1）。这证明 $F$ 在 $x$ 的某\Term{neighborhood}{邻域}中\Term{locally-constant}{局部常值}（在此甚至为\Term{constant-sheaf}{常层}）。证毕。
\end{proof}
\begin{SGAresult}{定义}{2.12}
若对 $X$ 的任意一对\Term{geometric-point}{几何点} $P,Q$，都存在\Term{geometric-point}{几何点}
\[
P=P_0,\ldots,P_n;\qquad Q_1,\ldots,Q_n=Q
\]
及\Term{specialization}{特殊化}（第 VIII 讲 7.2）$P_i\to Q_i$ 和 $P_{i-1}\to Q_i$（$i=1,\ldots,n$），就称\Term{scheme}{概形} $X$ \Term{arcwise-connected}{道路连通}。若对每个\Term{etale}{平展}的 $U\to X$，都存在 $U$ 的一个\Term{etale}{平展}\Term{covering}{覆盖} $\{U_i\to U\}$，使每个 $U_i$ 都\Term{arcwise-connected}{道路连通}，就称 $X$（关于\Term{etale}{平展}\Term{topology}{拓扑}）\Term{locally-arcwise-connected}{局部道路连通}。\footnote{这一术语由 S. Lubkin 提出。}
\end{SGAresult}
\SourcePage{22}
立即可验证，\Term{locally-noetherian}{局部诺特}\Term{prescheme}{预概形}是\Term{locally-arcwise-connected}{局部道路连通}的。
\begin{SGAresult}{命题}{2.13}
(i) 设 $X$ \Term{locally-arcwise-connected}{局部道路连通}，$F$ 是 $X$ 上的\Term{set}{集合}\Term{sheaf}{层}（分别为群\Term{sheaf}{层}、$A$-模\Term{sheaf}{层}）。设 $F$ 的\Term{fiber}{纤维}\Term{finite}{有限}（分别为\Term{finite}{有限}、\Term{finite-presentation}{有限表现}）。则 $F$ \Term{locally-constant}{局部常值}，当且仅当对\Term{geometric-point}{几何点}的每个\Term{specialization}{特殊化} $P\to Q$，\Term{specialization}{特殊化}\Term{morphism}{态射}（第 VIII 讲 7.7）$F_Q\to F_P$ 都是\Term{bijection}{双射}。

(ii) 设 $X$ \Term{locally-noetherian}{局部诺特}，且 $F$ 的\Term{fiber}{纤维}\Term{finite}{有限}（分别为\Term{finite}{有限}、\Term{finite-presentation}{有限表现}）。若对每个\Term{specialization}{特殊化} $P\to Q$，\Term{specialization}{特殊化}\Term{morphism}{态射} $F_Q\to F_P$ 都是\Term{injectivity}{单射}，则 $F$ \Term{constructible}{可构造}。

(iii) 设 $X$ \Term{locally-noetherian}{局部诺特}，\Term{set}{集合}\Term{sheaf}{层} $F$ 的\Term{fiber}{纤维}\Term{finite}{有限}。令 $c:X\to\mathbb N$ 为将 $x\in X$ 映到 $F$ 在 $x$ 上方某\Term{geometric-point}{几何点} $\bar x$ 处的\Term{fiber}{纤维} $F_{\bar x}$ 的元素个数的函数。则 $F$ \Term{constructible}{可构造}，当且仅当 $c$ 是\Term{constructible}{可构造}函数，即对每个 $n\in\mathbb N$，$c^{-1}(n)$ 都\Term{constructible}{可构造}。
\end{SGAresult}
\begin{proof}
(i) 设 $P$ 为 $X$ 的\Term{geometric-point}{几何点}。我们要找到 $P$ 的一个\Term{etale}{平展}\Term{neighborhood}{邻域}，即带 $P$ 标点的\Term{etale}{平展} $U\to X$，使 $F$ 在 $U$ 上成为\Term{constant-sheaf}{常层}。设 $V_0$ 为 $P$ 的一个\Term{etale}{平展}\Term{neighborhood}{邻域}，使 $F(V_0)\to F_P$ \Term{surjectivity}{满射}。由\Term{finiteness}{有限性}假设，这样的 $V_0$ 存在。取 $P$ 在 $V_0$ 中的一个\Term{etale}{平展}\Term{neighborhood}{邻域} $U\to V_0$，使 $U$ \Term{arcwise-connected}{道路连通}。则 $F(U)\to F_P$ 仍为\Term{surjectivity}{满射}。取\Term{subset}{子集} $S\subset F(U)$，使它\Term{bijection}{双射}到 $F_P$。若 $F$ 为群\Term{sheaf}{层}（分别为 $A$-模\Term{sheaf}{层}），容易看出必要时改变 $U$ 后，可使该\Term{set}{集合}还是 $F(U)$ 的一个\Term{subgroup}{子群}（分别为\Term{submodule}{子模}）。由此得到\Term{morphism}{态射} $S_U\to F|U$，其中 $S_U$ 为值为 $S$ 的\Term{constant-sheaf}{常层}。有
\[
(S_U)_P\xrightarrow{\sim}F_P.
\]
由于 $F$（分别为 $S_U$）的\Term{fiber}{纤维}之间每个\Term{specialization}{特殊化}\Term{morphism}{态射}都为\Term{bijection}{双射}，且 $U$ \Term{arcwise-connected}{道路连通}，得到 $S_U\xrightarrow{\sim}F|U$，结论即得。

(ii) 断言是\Term{local-scheme}{局部}的，故可设 $X$ \Term{noetherian}{诺特}。设 $P$ 为 $X$ 某个\Term{maximal-point}{极大点}上方的\Term{geometric-point}{几何点}。存在 $P$ 的一个\Term{irreducible}{不可约}的、因而\Term{arcwise-connected}{道路连通}的\Term{etale}{平展}\Term{neighborhood}{邻域} $V\to X$，使 $F(V)$ \Term{surjectivity}{满射}到 $F_P$。于是它对 $V$ 的每个\Term{geometric-point}{几何点} $Q$ 都\Term{surjectivity}{满射}到 $F_Q$；\Term{equivalence}{等价}地，对 $V$ 在 $X$ 中的像 $U$ 的每个\Term{geometric-point}{几何点}都如此，因为这样的点是 $P$ 的\Term{specialization}{特殊化}，且 $F_Q\to F_P$ 为\Term{injectivity}{单射}。所以 $U$ 中每个\Term{specialization}{特殊化}\Term{morphism}{态射}都为\Term{bijection}{双射}，故 $F|U$ \Term{locally-constant}{局部常值}。由\Term{noetherian-induction}{诺特归纳}法可设 $F|(X-U)$ \Term{constructible}{可构造}，结论即得。

(iii) 显然只须处理函数 $c$ 为常值的情形。设 $V\to X$ 是 $X$ 某个\Term{maximal-point}{极大点}上方的\Term{geometric-point}{几何点} $P$ 的一个\Term{etale}{平展}\Term{neighborhood}{邻域}，使 $F(V)$ \Term{surjectivity}{满射}到 $F_P$。取 $F(V)$ 的\Term{subset}{子集} $S$，使 $S\xrightarrow{\sim}F_P$。
\SourcePage{23}
则对\Term{constant-sheaf}{常层} $S_V$，以及\Term{geometric-point}{几何点}的每个\Term{specialization}{特殊化} $P\to Q$，有\Term{commutative-diagram}{交换图}
\[
\begin{tikzcd}
(S_V)_P\arrow[r,"\sim"]&F_P\\
(S_V)_Q\arrow[u,"\sim"]\arrow[r]&F_Q\arrow[u].
\end{tikzcd}
\]
因为 $c(F_Q)=c(F_P)$，\Term{arrow}{箭头} $F_Q\to F_P$ 为\Term{bijection}{双射}，所以由 (i)，$F$ 在 $P$ 的一个\Term{neighborhood}{邻域}内\Term{locally-constant}{局部常值}。再用\Term{noetherian-induction}{诺特归纳}法即得。
\end{proof}
下面的 (ii) 将在以后发挥技术\Term{action}{作用}，使某些验证可以归结到\Term{constant-sheaf}{常层}的情形。
\begin{SGAresult}{命题}{2.14}
设 $X$ 为\Term{scheme}{概形}。

(i) 对每个\Term{finite}{有限}且\Term{finite-presentation}{有限表现}的\Term{morphism}{态射} $f:X'\to X$，以及 $X'$ 上的每个\Term{constructible}{可构造}\Term{set}{集合}\Term{sheaf}{层}（分别为\Term{ind-finite-group}{ind-有限群}\Term{sheaf}{层}、$A$-模\Term{sheaf}{层}）$F$，$f_*(F)$ 都\Term{constructible}{可构造}。

(ii) 设 $X$ \Term{quasi-compact}{拟紧}且\Term{quasi-separated}{拟分离}，$F$ 是 $X$ 上的\Term{constructible}{可构造}\Term{set}{集合}\Term{sheaf}{层}（分别为……）。则存在一个\Term{finite}{有限}的\Term{finite-morphism}{有限态射}族 $(p_i:X'_i\to X_i)$，且对每个 $i$，存在 $X'_i$ 上的一个\Term{constructible}{可构造}\Term{constant-sheaf}{常集合层}（分别为……）$C_i$，以及一个\Term{monomorphism}{单态射}
\[
F\lhook\joinrel\longrightarrow\prod_ip_{i*}(C_i).
\]
\end{SGAresult}
\begin{proof}
先证明 (i)。按定义可设 $X$ \Term{affine}{仿射}；再利用 2.7.4，标准的取\Term{limit}{极限}方法将我们归结到 $X$ \Term{noetherian}{诺特}的情形。再次取\Term{limit}{极限}，仍利用 2.7.4，并利用 2.4 (v)，可归结到 $X$ 为某个域的谱的情形。此时结论归结到 $f_*(F)$ 在 $X$ 上某\Term{geometric-point}{几何点}处\Term{fiber}{纤维}的相应断言，而这由第 VIII 讲 5.5、5.8 立即得到。

证明 (ii)。由 2.4 (i)，存在将 $X$ 分为子\Term{scheme}{概形} $Y_i$ 的分划，其\Term{inclusion-morphism}{包含态射} $u_i:Y_i\to X$ 为\Term{finite-presentation}{有限表现}，使 $F$ 在每个 $Y_i$ 上的\Term{restriction}{限制} $F_i$ 为\Term{locally-constant}{局部常值}\Term{sheaf}{层}。立即可知，\Term{canonical-homomorphism}{典范同态} $F\to u_{i*}(F_i)$ 定义一个单\Term{homomorphism}{同态} $F\to\prod_iu_{i*}(F_i)$。

在上述构造中，甚至可使每个 $F_i$ 不仅\Term{locally-constant}{局部常值}，而且\Term{isotrivial}{等平凡}（isotrivial），即存在一个\Term{finite}{有限}\Term{etale}{平展}\Term{surjectivity}{满射} $q_i:Y'_i\to Y_i$，使 $q_i^*(F_i)$ 是 $Y'_i$ 上的\Term{constant-sheaf}{常层}，记为 $C_{i,Y'_i}$。例如，容易由定义、2.8.1 和第 VIII 讲 1.1 得到这一点。另一方面，由“主定理”（Main Theorem）的 EGA IV 8.12.6 \Term{form}{形式}，
\SourcePage{24}
存在一个\Term{finite-morphism}{有限态射} $p_i:Z_i\to X$ 以及 $X$-\Term{scheme}{概形}的一个\Term{open-immersion}{开浸入} $v_i:Y'_i\to Z_i$。

由于 $F_i$ \Term{embedding}{嵌入} $q_{i*}(q_i^*(F_i))=q_{i*}(C_{i,Y'_i})$，$F$ \Term{embedding}{嵌入}各 $u_{i*}(q_{i*}(C_{i,Y'_i}))=p_{i*}(v_{i*}(C_{i,Y'_i}))$ 的积。若每个 $Z_i$ 都\Term{normal}{正规}，则下面的引理 2.14.1 意味着 $v_{i*}(C_{i,Y'_i})\simeq C_{i,Z_i}$，所以 $F$ \Term{embedding}{嵌入}各 $p_{i*}(C_{i,Z_i})$ 的积，证明即告完成。

一般地，引入 $Z_i$ 的\Term{normalization}{正规化} $\overline Z_i$、$Y'_i$ 在 $\overline Z_i$ 中的\Term{inverse-image}{逆像} $\overline Y'_i$，以及\Term{immersion}{浸入} $\overline v_i:\overline Y'_i\to\overline Z_i$。立即可见 $\overline Y'_i$ 包含 $\overline Z_i$ 的各\Term{maximal-point}{极大点}，所以在 $\overline Z_i$ 中\Term{dense}{稠密}。设 $r_i:\overline Z_i\to Z_i$ 和 $s_i:\overline Y'_i\to Y'_i$ 为\Term{projective-module}{投射}。由于后者\Term{surjectivity}{满射}，$C_{i,Y'_i}$ \Term{embedding}{嵌入} $s_{i*}(C_{i,\overline Y'_i})$，所以 $v_{i*}(C_{i,Y'_i})$ \Term{embedding}{嵌入}
\[
v_{i*}(s_{i*}(C_{i,\overline Y'_i}))
 =r_{i*}(\overline v_{i*}(C_{i,\overline Y'_i})),
\]
由下面的 2.14.1，后者\Term{isomorphism}{同构}于 $r_{i*}(C_{i,\overline Z_i})$。因此 $F$ \Term{embedding}{嵌入}各
\[
p_{i*}(r_{i*}(C_{i,\overline Z_i}))=t_{i*}(C_{i,\overline Z_i})
\]
的积，其中 $t_i=p_ir_i$ 为\Term{canonical}{典范}\Term{projective-module}{投射} $\overline Z_i\to X$。若对每个 $i$，此\Term{morphism}{态射}均\Term{finite}{有限}，则已证明完毕。一般只能说 $\overline Z_i$ 在 $X$ 上整，因而形如 $\operatorname{Spec}(\mathcal A_i)$，其中 $\mathcal A_i$ 是一个\Term{quasi-coherent}{拟凝聚}的 $\mathcal O_X$-\Term{algebra}{代数}\Term{sheaf}{层}，且在 $\mathcal O_X$ 上整。由 EGA IV 1.7.7，$\mathcal A_i$ 是其在 $X$ 上\Term{finite}{有限}的子\Term{algebra}{代数} $\mathcal A_{i,j}$ 的\Term{inductive-limit}{归纳极限}，所以 $\overline Z_i$ 是在 $X$ 上\Term{finite}{有限}的\Term{prescheme}{预概形} $Z_{ij}=\operatorname{Spec}(\mathcal A_{ij})$ 的\Term{projective-limit}{投射极限}（第 VII 讲 5）。应用第 VII 讲 5.11，$t_{i*}(C_{i,\overline Z_i})$ 是各 $t_{ij*}(C_{i,Z_{ij}})$ 的\Term{inductive-limit}{归纳极限}，其中 $t_{ij}:Z_{ij}\to X$ 是\Term{canonical}{典范}\Term{projective-module}{投射}。由 2.7.3，\Term{homomorphism}{同态} $F\to t_{i*}(C_{i,\overline Z_i})$ 因而可经过某个\Term{homomorphism}{同态} $F\to t_{ij*}(C_{i,Z_{ij}})$ \Term{resolution}{分解}，这里 $j=j(i)$ 选取适当。现在置 $X'_i=Z_{i,j(i)}$，便得到 2.14 所述的\Term{embedding}{嵌入}。还须证明下述引理。
\end{proof}
\begin{SGAresult}{引理}{2.14.1}
设 $X$ 为\Term{geometrically-unibranch}{几何单分支}\Term{scheme}{概形}（例如\Term{normal}{正规}\Term{scheme}{概形}），$i:U\to X$ 为一个\Term{quasi-compact}{拟紧}\Term{open-immersion}{开浸入}。则 $U$ 的\Term{closure}{闭包} $Z$ 赋予诱导的\Term{reduced}{既约}结构后仍为\Term{geometrically-unibranch}{几何单分支}的；更精确地，对每个 $z\in Z$，有
\[
\mathcal O_{Z,z}\xleftarrow{\sim}\mathcal O_{X_{\mathrm{red}},z}.
\]
另设 $i$ 为\Term{dominant-morphism}{支配态射}，$C$ 为一个\Term{set}{集合}，$C_U$ 为 $U$ 上由 $C$ 定义的\Term{constant-sheaf}{常层}。则自然\Term{homomorphism}{同态} $C_X\to i_*(C_U)$ 为\Term{isomorphism}{同构}。将\Term{open-immersion}{开浸入} $i$ 换成一个包含 $i:x\to X$，其中 $x$ 为 $X$ 的\Term{maximal-point}{极大点}，仍有同样的结论。
\end{SGAresult}
\begin{proof}
由 EGA IV 2.3.11，第一条断言归结到 $X$ \Term{local-scheme}{局部}且\Term{closed-point}{闭点}为 $z$ 的情形；此时 $X_{\mathrm{red}}$ 整且 $U$ \Term{nonempty}{非空}，所以 $Z=X_{\mathrm{red}}$，断言显然成立（因此在这里，与其假设 $X$ \Term{geometrically-unibranch}{几何单分支}，只须假设 $X_{\mathrm{red}}$ 的各\Term{local-ring}{局部环}整）。对于第二条断言，逐\Term{fiber}{纤维}考察，利用第 VIII 讲 5.3 的同一计算，可归结到 $X$ \Term{strictly-local}{严格局部}的情形，这时只须证明 $H^0(X,C_X)\to H^0(U,C_U)$ 为\Term{bijection}{双射}。此时 $X$ 仍\Term{irreducible}{不可约}，而 $U$ 为其\Term{nonempty}{非空}\Term{open-subset}{开集}，所以也\Term{irreducible}{不可约}；上述映射的两端都等同于 $C$，所需结论即得。
\SourcePage{25}
（这里使用 EGA IV 18.8.14 \Term{form}{形式}的“\Term{geometrically-unibranch}{几何单分支}”概念，即 $X_{\mathrm{red}}$ 的各\Term{strict-localization}{严格局部化}环都是\Term{integral-domain}{整环}。）\Term{maximal-point}{极大点}的包含情形完全同样处理（为方便以后引用，在此一并陈述）。
\end{proof}
\textbf{注记 2.14.2。} 当 $X$ \Term{noetherian}{诺特}时，上述证明（或它的一个直接推论）表明，在 2.14 (ii) 中可设 $X'_i$ 整；若进一步假设 $X$ 为\Term{universal}{普遍}日本的，即其\Term{affine-open}{仿射开集}的环都是\Term{universal}{普遍}日本环（EGA $0_{\mathrm{IV}}$ 23.1.1、IV 7.7.2），则还可设 $X'_i$ \Term{normal}{正规}。
\SGAsection{3}{Kummer 理论与 Artin--Schreier 理论}
\textbf{3.0.} 用 $(G_m)_X$ 表示 $X$ 上\Term{multiplicative-group}{乘法群} $\operatorname{Spec}\mathcal O_X[t,t^{-1}]$ 所表示的\Term{sheaf}{层}。因此，对\Term{etale}{平展}的 $U\to X$，$(G_m)_X$ 在 $U$ 上的\Term{section}{截面}是 $\Gamma(U,\mathcal O_U)$ 的\Term{invertible-element}{可逆元}素 $\Gamma(U,\mathcal O_U^*)$。设 $n\in\mathbb N$。$(G_m)_X$ 中 $n$ 次幂映射的核为“$n$ 次\Term{root-of-unity}{单位根}\Term{sheaf}{层}”，记为 $(\mu_n)_X$。于是有\Term{exact-sequence}{正合列}
\[
0\longrightarrow(\mu_n)_X\longrightarrow(G_m)_X\xrightarrow{n}(G_m)_X.
\tag{3.1}
\]
若 $n$ 在 $X$ 上\Term{invertible}{可逆}（即 $n$ 与 $X$ 的各\Term{residue-characteristic}{剩余特征}互素），则 $n$ 次幂映射作为\Term{sheaf}{层}\Term{morphism}{态射}\Term{surjectivity}{满射}：
\begin{SGAresult}{Kummer 理论}{3.2}
若 $n\in\mathbb N$ 在 $X$ 上\Term{invertible}{可逆}，则有\Term{exact-sequence}{正合列}
\[
0\longrightarrow(\mu_n)_X\longrightarrow(G_m)_X\xrightarrow{n}(G_m)_X\longrightarrow0.
\]
\end{SGAresult}
\begin{proof}
设 $u\in(G_m)_X(U)$，$U\to X$ \Term{etale}{平展}。须找到 $U$ 的一个\Term{etale}{平展}\Term{covering}{覆盖} $\{U_i\to U\}$，使 $u$ 在每个 $U_i$ 上成为一个 $n$ 次幂。由于 $n$ 在 $U$ 上\Term{invertible}{可逆}，方程
\[
T^n-u=0
\]
在 $U$ 上“\Term{separable}{可分}”，即
\[
U'=\operatorname{Spec}\bigl(\mathcal O_U[T]/(T^n-u)\bigr)
\]
在 $U$ 上\Term{etale}{平展}。$U'\to U$ \Term{surjectivity}{满射}，故是一个\Term{covering}{覆盖}；而 $u$ 在 $U'$ 上成为一个 $n$ 次幂，证明完毕。
\end{proof}
指出一点：当 $n$ 在 $X$ 上\Term{invertible}{可逆}时，$(\mu_n)_X$ 为\Term{locally-constant}{局部常值}\Term{sheaf}{层}，实际上\Term{local-isomorphism}{局部同构}于\Term{sheaf}{层} $(\mathbb Z/n\mathbb Z)_X$，而且在一些重要情形中它是\Term{constant-sheaf}{常层}。因此 Kummer 理论给出以某些\Term{constant-sheaf}{常层}为值的 $X$ 的\Term{cohomology}{上同调}的信息。事实上，按定义
\[
H^0(X,(G_m)_X)=\Gamma(X,\mathcal O_X^*),
\]
而在\Term{dimension}{维数} $1$ 有下述定理。
\SourcePage{26}
\begin{SGAresult}{定理}{3.3}
（“Hilbert 第 90 定理”）有一个\Term{isomorphism}{同构}
\[
H^1(X,(G_m)_X)\xrightarrow{\sim}\operatorname{Pic}X,
\]
其中 $\operatorname{Pic}X$ 为 $X$ 上\Term{invertible}{可逆}\Term{sheaf}{层}的类群。
\end{SGAresult}
\begin{proof}
这就是说，\Term{canonical}{典范}\Term{morphism}{态射}
\[
H^1(X_{\mathrm{zar}},(G_m)_X)\longrightarrow H^1(X_{\text{ét}},(G_m)_X)
\]
为\Term{bijection}{双射}；\Term{equivalence}{等价}地，
\[
R^1\varepsilon_*(G_m)_X=0,
\]
其中 $\varepsilon_*$ 是显然的\Term{site}{景}\Term{morphism}{态射} $\varepsilon:X_{\text{ét}}\to X_{\mathrm{zar}}$ 的\Term{direct-image}{直像}（第 VIII 讲 4）。于是归结到 $X$ \Term{affine}{仿射}的情形。由于可以用 \Term{cech-cohomology}{Čech 上同调}计算 $H^1$（第 V 讲 2.5），而对于\Term{quasi-compact}{拟紧}的 $X$，只须取\Term{quasi-compact}{拟紧}\Term{covering}{覆盖}，故这立即由\Term{sheaf}{层}的\Term{descent}{下降}理论得到（参看 FGA 190，或 SGA 第 VIII 讲 1）。
\end{proof}
\textbf{3.3.1.} 设 $X$ \Term{noetherian}{诺特}且无\Term{embedded-component}{嵌入分支}，$x_j$ 为 $X$ 的各\Term{maximal-point}{极大点}。令 $R_j$ 为 $X$ 在 $x_j$ 处的\Term{artinian-local-ring}{阿廷局部环}，$i_j:\operatorname{Spec}R_j\to X$ 为\Term{inclusion-morphism}{包含态射}。有 $X_{\text{ét}}$（分别为 $X_{\mathrm{zar}}$）上\Term{sheaf}{层}的\Term{canonical}{典范}\Term{injectivity}{单射}
\[
(G_m)_X\longrightarrow\prod_ji_{j*}(G_m)_{\operatorname{Spec}R_j}.
\]
其\Term{cokernel}{余核} $D$ 称为 $X_{\text{ét}}$（分别为 $X_{\mathrm{zar}}$）上的\Term{cartier-divisor}{Cartier 除子}\Term{sheaf}{层}。\footnote{在 EGA IV 21 中，以“\Term{divisor}{除子}”简称“\Term{cartier-divisor}{Cartier 除子}”。}
\begin{SGAresult}{推论}{3.4}
设 $X$ \Term{noetherian}{诺特}且无\Term{embedded-component}{嵌入分支}，$\varepsilon:X_{\text{ét}}\to X_{\mathrm{zar}}$ 为\Term{canonical}{典范}\Term{site}{景}\Term{morphism}{态射}，$D$ 为 $X_{\text{ét}}$ 上的\Term{cartier-divisor}{Cartier 除子}\Term{sheaf}{层}。则 $\varepsilon_*D$ 为 $X_{\mathrm{zar}}$ 上的\Term{cartier-divisor}{Cartier 除子}\Term{sheaf}{层} $D_{\mathrm{zar}}$；特别地，有
\[
H^0(X_{\text{ét}},D)\simeq H^0(X_{\mathrm{zar}},D_{\mathrm{zar}}).
\]
\end{SGAresult}
\begin{proof}
对\Term{exact-sequence}{正合列}
\[
0\longrightarrow(G_m)_X\longrightarrow\prod_ji_{j*}(G_m)_{\operatorname{Spec}R_j}\longrightarrow D\longrightarrow0
\]
应用 $\varepsilon_*$，由于 $R^1\varepsilon_*(G_m)_X=0$（3.3），仍得到\Term{exact-sequence}{正合列}，故第一条断言成立。
\end{proof}
若 $X$ 有\Term{characteristic}{特征} $p\ne0$，在研究 $p$-\Term{torsion-coefficient}{挠系数}时，下面的结果有时将代替 3.2。
\begin{SGAresult}{Artin--Schreier 理论}{3.5}
若 $X$ 的\Term{characteristic}{特征}为 $p>0$，则有\Term{exact-sequence}{正合列}
\[
0\longrightarrow(\mathbb Z/p\mathbb Z)_X\longrightarrow\mathcal O_X\xrightarrow{\wp}\mathcal O_X\longrightarrow0,
\]
\SourcePage{27}
其中 $\wp$ 为由 $\wp(f)=f^p-f$ 定义的\Term{additive-group}{加法群}\Term{morphism}{态射}。
\end{SGAresult}
\begin{proof}
$\wp$ 的核是 $\mathcal O_X$ 中\Term{local-scheme}{局部}属于\Term{prime-field}{素域}的\Term{section}{截面}所成的\Term{sheaf}{层}，所以是\Term{constant-sheaf}{常层}。$\wp$ 在\Term{etale}{平展}\Term{topology}{拓扑}中\Term{surjectivity}{满射}，因为方程 $T^p-T-a$（$a\in\Gamma(U,\mathcal O_U)$）在\Term{characteristic}{特征} $p$ 时总是“\Term{separable}{可分}”的，即定义 $U$ 的一个\Term{etale-covering}{平展覆盖}。应用此\Term{exact-sequence}{正合列}时，当然还要使用第 VII 讲 4.2。
\end{proof}
\begin{SGAresult}{命题}{3.6}
(i) 设 $X$ 为\Term{scheme}{概形}，$i:x\to X$ 为一个点的包含。则
\[
H^1(X,i_*(\mathbb Z)_x)=0,
\]
其中 $(\mathbb Z)_x$ 为值为 $\mathbb Z$ 的“\Term{constant-sheaf}{常层}”（参看 2.0）。

(ii) 若 $X$ \Term{irreducible}{不可约}，且对 $X$ 的每个\Term{geometric-point}{几何点} $\bar x$，$X$ 在 $\bar x$ 处的\Term{strict-localization}{严格局部化} $X_{\bar x}$ \Term{irreducible}{不可约}（应当说“$X$ \Term{geometrically-unibranch}{几何单分支}”，参看 EGA IV 18.8.14），则有
\[
H^1(X,\mathbb Z_X)=0.
\]
\end{SGAresult}
\begin{proof}
(i) 由 \Term{leray-spectral-sequence}{Leray 谱序列}
\[
E_2^{pq}=H^p(X,R^qi_*(\mathbb Z)_x)\Longrightarrow H^{p+q}(x,(\mathbb Z)_x)
\]
得到一个\Term{injectivity}{单射} $H^1(X,i_*(\mathbb Z)_x)\to H^1(x,(\mathbb Z)_x)$。由于 $x=\operatorname{Spec}k(x)$ 是域的谱，众所周知 $H^1(x,(\mathbb Z)_x)=0$（这由 $H^q(x,F)$ 在 $q>0$ 时总是\Term{torsion-group}{挠群}这一事实——参看第 VII 讲 2.2 和 CG——以及\Term{exact-sequence}{正合列} $0\to\mathbb Z\xrightarrow{n}\mathbb Z\to\mathbb Z/n\to0$ 得到），结论即得。

(ii) 令 $i:x\to X$ 为 $X$ 的\Term{generic-point}{泛点}的包含。立即可见，在 (ii) 的条件下，\Term{canonical}{典范}\Term{morphism}{态射} $(\mathbb Z)_X\to i_*(\mathbb Z)_x$ 为\Term{bijection}{双射}（反过来也成立！），于是 (ii) 归结为 (i)。
\end{proof}
\textbf{注记 3.7。} 在 3.6 (i)、(ii) 中，可将 $\mathbb Z$ 换成任意\Term{torsion-free}{无挠}\Term{abelian-group}{阿贝尔群}。另一方面，若在 (ii) 中去掉 $X$ \Term{geometrically-unibranch}{几何单分支}的假设，结论就不再成立。例如取 $X$ 为域 $K$ 上一条有\Term{ordinary-double-point}{普通二重点}的\Term{algebraic-curve}{代数曲线}即可看出（请在此情形计算 $H^1(X,\mathbb Z)$！）。

\SGAsection{4}{代数曲线的情形}
\textbf{4.0.} 设 $X$ 为\Term{dimension}{维数} $1$ 的\Term{noetherian}{诺特}\Term{scheme}{概形}。令 $R(X)$ 为 $X$ 上的\Term{rational-function-ring}{有理函数环}，$i:\operatorname{Spec}R(X)\to X$ 为\Term{inclusion-morphism}{包含态射}。环 $R(X)$ 是\Term{artinian-ring}{阿廷环}，等于 $X$ 在各\Term{maximal-point}{极大点}处的\Term{local-ring}{局部环}之积，可以应用第 VIII 讲 2 的结果。
\SourcePage{28}
设 $F$ 为 $\operatorname{Spec}R(X)$ 上的\Term{abelian-sheaf}{阿贝尔层}，考虑 $X$ 上的\Term{sheaf}{层} $R^qi_*F$，$q>0$。由第 VIII 讲 5.2 立即得到，$R^qi_*F$ 在 $X$ 的\Term{maximal-point}{极大点}上方任一\Term{geometric-point}{几何点}处的\Term{fiber}{纤维}为零。由于 $X$ 的\Term{dimension}{维数}为 $1$，这样的\Term{sheaf}{层}具有很特殊的\Term{form}{形式}：
\begin{SGAresult}{引理}{4.1}
设 $X$ \Term{noetherian}{诺特}，$F$ 为 $X$ 上的\Term{sheaf}{层}。以下条件\Term{equivalence}{等价}（满足这些条件的\Term{sheaf}{层}称为“\Term{skyscraper-sheaf}{摩天层}”，skyscraper sheaf）：

(i) 对 $X$ 的每个非\Term{closed-point}{闭点} $x\in X$ 上方的\Term{geometric-point}{几何点} $\bar x$，$F$ 在 $\bar x$ 处的\Term{fiber}{纤维}为零。

(ii) 每个\Term{section}{截面} $s\in F(X')$（$X'\to X$ \Term{etale}{平展}且为\Term{finite-type}{有限型}）除在 $X$ 的\Term{finite}{有限}个\Term{closed-point}{闭点}外均为零。

(iii) 有 $F\simeq\bigoplus_xi_{x*}i_x^*F$，其中 $x$ 遍历 $X$ 的所有\Term{closed-point}{闭点}，$i_x:x\to X$ 为包含。
\end{SGAresult}
\begin{proof}
设 (i) 成立，取 $z\in F(X')$（$X'\to X$ \Term{etale}{平展}且为\Term{finite-type}{有限型}）。则 $z$ 的\Term{support}{支集}（第 VIII 讲 6.6）是 $X'$ 中仅由\Term{closed-point}{闭点}组成的闭集，因此\Term{finite}{有限}（因为 $X'$ \Term{noetherian}{诺特}），故 (i) 推出 (ii)。(ii) 推出 (i) 由\Term{fiber}{纤维}的计算（第 VIII 讲 3.9）得到；(iii) 推出 (i)，因为对 $X$ 的\Term{closed-point}{闭点} $x$ 上的任意\Term{sheaf}{层}族 $(G_x)$，$X$ 上的\Term{sheaf}{层} $F=\bigoplus_xi_{x*}(G_x)$ 都满足 (i)：这是因为\Term{fibre-functor}{纤维函子}与\Term{direct-sum}{直和}交换，而 $i_{x*}(G_x)$ 的\Term{support}{支集}显然包含在 $\{x\}$ 中。

还须证明 (ii) 推出 (iii)。考虑\Term{canonical}{典范}\Term{morphism}{态射}
\[
F\longrightarrow\prod_xi_{x*}i_x^*F,
\]
其中 $x$ 遍历 $X$ 的所有\Term{closed-point}{闭点}。由 (ii)，它经过
\[
F\longrightarrow\bigoplus_xi_{x*}i_x^*F
\]
\Term{resolution}{分解}，便得到所需\Term{morphism}{态射}。它为\Term{bijection}{双射}的验证留给读者。
\end{proof}
\textbf{4.1.1.} 对上面引入的\Term{sheaf}{层} $R^qi_*F$ 应用此引理，有
\[
R^qi_*F\xrightarrow{\sim}\bigoplus_xi_{x*}i_x^*(R^qi_*F),
\]
其中 $x$ 遍历 $X$ 的所有\Term{closed-point}{闭点}。由第 VIII 讲 5，$R^qi_*F$ 在 $X$ 的\Term{geometric-point}{几何点} $\bar x$ 处的\Term{fiber}{纤维}为
\[
(R^qi_*F)_{\bar x}\simeq H^q(\operatorname{Spec}R(X_{\bar x}),F),
\]
其中 $R(X_{\bar x})\simeq R(X)\otimes_{\mathcal O_X}\mathcal O_{X_{\bar x}}$，$X_{\bar x}$ 为 $X$ 在 $\bar x$ 处的\Term{strict-localization}{严格局部化}（第 VIII 讲 4）；仍用 $F$ 表示 $F$ 在 $\operatorname{Spec}R(X_{\bar x})$ 上诱导的\Term{sheaf}{层}。
\SourcePage{29}
现设 $X$ 每个\Term{closed-point}{闭点} $x$ 的\Term{residue-field}{剩余域}都\Term{separably-closed}{可分闭}。于是有（也参看第 VI 讲 1）
\[
\begin{aligned}
H^p(X,R^qi_*F)
&\simeq H^p\left(X,\bigoplus_xi_{x*}i_x^*(R^qi_*F)\right)\\
&\simeq\bigoplus_xH^p(X,i_{x*}i_x^*(R^qi_*F))\\
&\simeq\bigoplus_xH^p(x,i_x^*(R^qi_*F))\\
&\simeq\begin{cases}
\displaystyle\bigoplus_xH^q(\operatorname{Spec}R(X_{\bar x}),F),&p=0,\\
0,&p>0.
\end{cases}
\end{aligned}
\]
由此得到：
\begin{SGAresult}{推论}{4.2}
设 $X$ 为\Term{dimension}{维数} $1$ 的\Term{noetherian}{诺特}\Term{scheme}{概形}，使对 $X$ 的每个\Term{closed-point}{闭点} $x$，\Term{residue-field}{剩余域} $k(x)$ 都\Term{separably-closed}{可分闭}。设 $F$ 为 $\operatorname{Spec}R(X)$ 上的\Term{abelian-sheaf}{阿贝尔层}（其中 $R(X)$ 为 $X$ 上\Term{rational-function-ring}{有理函数环}）。考虑 \Term{leray-spectral-sequence}{Leray 谱序列}（\SGAsourcegap{原文此处引用留空；补源原版第 IX 讲第 32 页亦同}）
\[
E_2^{pq}=H^p(X,R^qi_*F)\Longrightarrow H^{p+q}(\operatorname{Spec}R(X),F).
\]
有
\[
\begin{aligned}
E_2^{pq}&=0&&\text{若 }p,q>0,\\
E_2^{0q}&=\bigoplus_xH^q(\operatorname{Spec}R(X_{\bar x}),F)&&\text{若 }q>0,
\end{aligned}
\]
其中 $x$ 遍历 $X$ 的所有\Term{closed-point}{闭点}，$X_{\bar x}$ 为 $X$ 在 $x$ 处的\Term{strict-localization}{严格局部化}。

因此，该\Term{spectral-sequence}{谱序列}给出\Term{exact-sequence}{正合列}
\[
\begin{aligned}
0&\longrightarrow H^1(X,i_*F)\longrightarrow H^1(\operatorname{Spec}R(X),F)
 \longrightarrow\bigoplus_xH^1(\operatorname{Spec}R(X_{\bar x}),F)\\
&\longrightarrow H^2(X,i_*F)\longrightarrow H^2(\operatorname{Spec}R(X),F)
 \longrightarrow\bigoplus_xH^2(\operatorname{Spec}R(X_{\bar x}),F)\\
&\longrightarrow\cdots.
\end{aligned}
\]
\end{SGAresult}
\textbf{4.3.} 现进一步假设，对某个给定素数 $\ell$，$\operatorname{Spec}R(X)$ 的 $\ell$-\Term{cohomological-dimension}{上同调维数}至多为 $1$（即对 $\operatorname{Spec}R(X)$ 上每个\Term{p-torsion-sheaf}{$\ell$-挠层} $F$，都有 $H^q(\operatorname{Spec}R(X),F)=0$，$q>1$）。令 $K^1,\ldots,K^m$ 为\Term{artinian-ring}{阿廷环} $R(X)$ 的各\Term{residue-field}{剩余域}，即 $X$ 的\Term{maximal-point}{极大点}的\Term{residue-field}{剩余域}。结合第 VIII 讲 2 的对应，$\operatorname{Spec}R(X)$ 的 $\ell$-\Term{cohomological-dimension}{上同调维数}为
\[
\operatorname{cd}_\ell(\operatorname{Spec}R(X))
 =\sup_i\{\operatorname{cd}_\ell(G(\overline K^{\,i}/K^i))\},
\]
其中 $\operatorname{cd}_\ell(G)$ 为\Term{profinite-group}{投有限群} $G$ 的 $\ell$-\Term{cohomological-dimension}{上同调维数}（参看 CG）。
\SourcePage{30}
对 $X$ 的一个\Term{strict-localization}{严格局部化} $\overline X(x)$，$R(\overline X(x))$ 的每个\Term{residue-field}{剩余域} $K_x$ 显然都等同于某个 $K^i$ 的一个\Term{separable}{可分}（无限）\Term{group-extension}{扩张}，所以由 CG II 4.1，命题 10，有
\[
\operatorname{cd}_\ell(G(\overline K_x/K_x))
 \le\operatorname{cd}_\ell(G(\overline K^{\,i}/K^i)).
\]
因此，$\ell$-\Term{cohomological-dimension}{上同调维数} $\operatorname{cd}_\ell(\operatorname{Spec}R(\overline X(x)))$ 也不超过 $1$。

对\Term{p-torsion-sheaf}{$\ell$-挠层} $F$ 应用 4.2 的\Term{exact-sequence}{正合列}，得到
\[
\left\{\begin{aligned}
0&\longrightarrow H^1(X,i_*F)\longrightarrow H^1(\operatorname{Spec}R(X),F)\\
&\longrightarrow\bigoplus_xH^1(\operatorname{Spec}R(\overline X(x)),F)
 \longrightarrow H^2(X,i_*F)\longrightarrow0,\\
&H^q(X,i_*F)=0\qquad\text{对 }q>2.
\end{aligned}\right.
\tag{4.3.1}
\]
注意，若 $X$ 是\Term{separably-closed-field}{可分闭域}上的\Term{algebraic-curve}{代数曲线}，则满足 $\operatorname{cd}_\ell(\operatorname{Spec}R(X))\le1$ 这一假设（Tsen 定理 [6，定理 4.3]）。Ogg [2] 和 Šafarevič 在这一情形研究过类似的\Term{exact-sequence}{正合列}。

\textbf{4.4.} 设 $\ell$ 在 $X$ 上\Term{invertible}{可逆}，并取 $F=(G_m)_{R(X)}$。应用 3.3 及 CG 第 I 章命题 12，仍在 $\operatorname{cd}_\ell(\operatorname{Spec}R(X))\le1$ 的假设下，得到
\[
\left\{\begin{aligned}
H^1(\operatorname{Spec}R(X),G_m)&=0,\\
H^q(\operatorname{Spec}R(X),G_m)&\underset{\ell}{=}0\qquad(q\ge2).
\end{aligned}\right.
\tag{4.4.1}
\]
这里等号下的 $\ell$ 意味着，左端的群（无论如何都是\Term{torsion-group}{挠群}）没有 $\ell$-挠元。

将 $R(X)$ 换成 $R(\overline X(x))$，仍有相同结果。由于 $\operatorname{Spec}R(X)$ 上的\Term{sheaf}{层} $G_m$ 显然在 $\operatorname{Spec}R(\overline X(x))$ 上诱导\Term{sheaf}{层} $G_m$，由 4.2 得到
\[
\left\{\begin{aligned}
H^1(X,i_*(G_m)_{R(X)})&=0,\\
H^q(X,i_*(G_m)_{R(X)})&\underset{\ell}{=}0\qquad(q>1).
\end{aligned}\right.
\tag{4.4.2}
\]
对 $X$ 上任意\Term{sheaf}{层} $F$，由 4.1 可知，\Term{canonical-homomorphism}{典范同态} $F\to i_*i^*F$ 的核 $P$ 及\Term{cokernel}{余核} $Q$ 都是“\Term{skyscraper-sheaf}{摩天层}”；因假设 $X$ 的\Term{closed-point}{闭点}处的\Term{residue-field}{剩余域}\Term{separably-closed}{可分闭}，对 $i>0$ 有 $H^i(X,P)=H^i(X,Q)=0$。由\Term{cohomology}{上同调}\Term{exact-sequence}{正合列}立即推出，当 $i\ge2$ 时，
\[
H^i(X,F)\longrightarrow H^i(X,i_*i^*F)
\]
为\Term{isomorphism}{同构}。结合 (4.4.2)，得到关系
\[
H^i(X,G_m)\underset{\ell}{=}0\qquad\text{对 }i\ge2,
\tag{4.5}
\]
\SourcePage{31}
它在 $X$ 为\Term{dimension}{维数} $1$ 的\Term{noetherian}{诺特}\Term{prescheme}{预概形}、满足 $\operatorname{cd}_\ell(R(X))\le1$、且\Term{closed-point}{闭点}处\Term{residue-field}{剩余域}均\Term{separably-closed}{可分闭}时成立。结合 Kummer 理论（3.2）及 3.3，得到：
\begin{SGAresult}{定理}{4.6}
设 $X$ \Term{noetherian}{诺特}且\Term{dimension}{维数}为 $1$，$n$ 为 $X$ 上\Term{invertible}{可逆}的整数。设对 $n$ 的每个素因子 $\ell$，都有 $\operatorname{cd}_\ell(R(X))\le1$；且对 $X$ 的每个\Term{closed-point}{闭点} $x$，\Term{residue-field}{剩余域} $k(x)$ 都\Term{separably-closed}{可分闭}。则
\[
H^q(X,(\mu_n)_X)=0\qquad\text{若 }q>2,
\]
而 $q\le2$ 时的\Term{cohomology}{上同调}由以下\Term{exact-sequence}{正合列}给出：
\[
\begin{aligned}
0&\longrightarrow H^0(X,(\mu_n)_X)\longrightarrow\Gamma(X,\mathcal O_X^*)
 \xrightarrow{n}\Gamma(X,\mathcal O_X^*)\\
&\longrightarrow H^1(X,(\mu_n)_X)\longrightarrow\operatorname{Pic}X
 \xrightarrow{n}\operatorname{Pic}X\longrightarrow H^2(X,(\mu_n)_X)\longrightarrow0.
\end{aligned}
\]
\end{SGAresult}
\begin{SGAresult}{推论}{4.7}
设 $X$ 为\Term{separably-closed-field}{可分闭域} $k$ 上的一条\Term{complete}{完备}\Term{connected}{连通}曲线，$k$ 的\Term{characteristic}{特征}与 $n$ 互素。则
\[
\begin{aligned}
H^0(X,(\mu_n)_X)&\simeq\mu_n(k),\\
H^1(X,(\mu_n)_X)&\simeq{}_n(\operatorname{Pic}X)
 =\{\operatorname{Pic}X\text{ 中阶整除 }n\text{ 的元素}\},\\
H^2(X,(\mu_n)_X)&\simeq(\operatorname{Pic}X)/n\simeq(\mathbb Z/n\mathbb Z)^c,
\end{aligned}
\]
其中 $c$ 为 $X$ 的\Term{irreducible-component}{不可约分支}的\Term{set}{集合}。
\end{SGAresult}
对于最后一条断言，“回忆”这样的 $X/\operatorname{Spec}k$ 的“Picard \Term{scheme}{概形}”[TDTE V] 有一个\Term{resolution}{分解}
\[
0\longrightarrow A\longrightarrow\operatorname{Pic}_{X/k}\longrightarrow\mathbb Z^c\longrightarrow0,
\tag{4.8}
\]
其中 $A$ 为 $\operatorname{Spec}k$ 上的\Term{connected}{连通}\Term{algebraic-group}{代数群}。而这样的 $A$ 中“乘以 $n$”的映射是\Term{etale}{平展}的，所以对 $k$-值点\Term{surjectivity}{满射}。由此得到
\[
{}_n(\operatorname{Pic}X)\simeq{}_n(A(k)),\qquad
(\operatorname{Pic}X)/n\simeq(\mathbb Z/n\mathbb Z)^c.
\]
若 $X$ 在 $\operatorname{Spec}k$ 上\Term{smooth}{光滑}，使 Jacobian $A$ 为\Term{abelian-variety}{阿贝尔簇}，则群 ${}_n(\operatorname{Pic}X)$ \Term{isomorphism}{同构}于 $(\mathbb Z/n)^{2g}$，其中 $g$ 为 $X$ 的\Term{genus}{亏格}。当然，也有 $\mu_n(k)\simeq\mathbb Z/n\mathbb Z$；因此，各\Term{cohomology-group}{上同调群}作为 $\mathbb Z/n$-模的秩为 $1,2g,1$，这正是所预期的结果。
\SourcePage{32}
\SGAsection{5}{迹方法}
\textbf{5.1.} 设 $X$ 为\Term{scheme}{概形}，$f:X'\to X$ 为一个\Term{etale-covering}{平展覆盖}，$F$ 为 $X$ 上的\Term{sheaf}{层}。\Term{adjunction}{伴随}公式给出一个\Term{restriction}{限制}\Term{morphism}{态射}
\[
F\xrightarrow{\mathrm{res}}f_*f^*F.
\]
由于 $f$ \Term{finite}{有限}，有 $H^q(X,f_*f^*F)\simeq H^q(X',f^*F)$，从而得到仍称为\Term{restriction}{限制}的\Term{morphism}{态射}
\[
H^q(X,F)\xrightarrow{\mathrm{res}}H^q(X',f^*F),
\tag{5.1.1}
\]
它正是显然的\Term{morphism}{态射}。但当 $f$ \Term{etale}{平展}时，还有另一个\Term{morphism}{态射}，称为\Term{trace-morphism}{迹态射}：
\[
f_*f^*F\xrightarrow{\mathrm{tr}}F,
\tag{5.1.2}
\]
它在\Term{cohomology}{上同调}上给出\Term{morphism}{态射}
\[
H^q(X',f^*F)\xrightarrow{\mathrm{tr}}H^q(X,F).
\tag{5.1.3}
\]
\Term{trace-morphism}{迹态射}由下列两条性质刻画：

(i) 它关于\Term{etale}{平展}\Term{topology}{拓扑}是\Term{local-scheme}{局部}的。

(ii) 若 $X'$ 是 $X$ 的 $d$ 个副本的\Term{disjoint-coproduct}{不交和}，使 $f_*f^*F\simeq F^d$，则迹为求和映射 $F^d\to F$。

由于每个\Term{etale-covering}{平展覆盖}都是\Term{locally-constant}{局部常值}的，唯一性由 (i)、(ii) 立即得到。有了唯一性，再由通常的\Term{descent}{下降}论证即得存在性。

由 (ii) 推出下式：
\[
\mathrm{tr}\circ\mathrm{res}:F\longrightarrow F\quad
\text{为乘以 }f\text{ 的局部次数的映射}.
\tag{5.1.4}
\]
若次数为常数 $d$，则\Term{morphism}{态射} $\mathrm{tr}\circ\mathrm{res}:H^q(X,F)\to H^q(X,F)$ 为乘以 $d$ 的映射。
\begin{SGAresult}{推论}{5.2}
设 $f:X'\to X$ 为常次数 $d$ 的\Term{etale-covering}{平展覆盖}。若 $F$ 为 $X$ 上的 $r$-\Term{torsion-sheaf}{挠层}，且 $(r,d)=1$，则 $H^q(X',f^*F)=0$ 蕴含 $H^q(X,F)=0$。
\end{SGAresult}
事实上，乘以 $d$ 在 $F$ 上诱导一个\Term{isomorphism}{同构}，因而在\Term{cohomology}{上同调}上诱导一个\Term{isomorphism}{同构}。若后一个\Term{isomorphism}{同构}经过零\Term{resolution}{分解}，则 $H^q(X,F)=0$。

\textbf{5.2.1.} 现设 $i:U\to X$ 为\Term{immersion}{浸入}，$f:U'\to U$ 为次数 $d$ 的\Term{etale-covering}{平展覆盖}，并设 $f$ 是某个\Term{finite-morphism}{有限态射} $g:X'\to X$ 经\Term{base-change}{基变换}诱导的，
\SourcePage{33}
即有\Term{cartesian-diagram}{笛卡儿图}
\[
\begin{tikzcd}
U'\arrow[r,"i'"]\arrow[d,"f"']&X'\arrow[d,"g"]\\
U\arrow[r,"i"]&X.
\end{tikzcd}
\]
设 $F$ 为 $U$ 上的\Term{sheaf}{层}。有
\[
\begin{tikzcd}
F\arrow[r,"\mathrm{res}"]\arrow[rr,bend right=20,"d"']&f_*f^*F\arrow[r,"\mathrm{tr}"]&F.
\end{tikzcd}
\]
因为 $i_*f_*=g_*i'_*$，有
\[
\begin{tikzcd}
i_*F\arrow[r,"\mathrm{res}"]\arrow[rr,bend right=20,"d"']&g_*i'_*f^*F\arrow[r,"\mathrm{tr}"]&i_*F.
\end{tikzcd}
\tag{5.3}
\]
又因 $i_!f_*\simeq g_*i'_!$（例如由第 VIII 讲 5.5 得到），有
\[
\begin{tikzcd}
i_!F\arrow[r,"\mathrm{res}"]\arrow[rr,bend right=20,"d"']&g_*i'_!f^*F\arrow[r,"\mathrm{tr}"]&i_!F.
\end{tikzcd}
\tag{5.4}
\]
从而
\[
\begin{tikzcd}[column sep=small]
H^q(X,i_*F)\arrow[r,"\mathrm{res}"]\arrow[rr,bend right=20,"d'"']&H^q(X',i'_*f^*F)\arrow[r,"\mathrm{tr}"]&H^q(X,i_*F)
\end{tikzcd}
\tag{5.3'}
\]
以及
\[
\begin{tikzcd}[column sep=small]
H^q(X,i_!F)\arrow[r,"\mathrm{res}"]\arrow[rr,bend right=20,"d"']&H^q(X',i'_!f^*F)\arrow[r,"\mathrm{tr}"]&H^q(X,i_!F).
\end{tikzcd}
\tag{5.4'}
\]
由 (5.4) 得到以下推论：
\begin{SGAresult}{命题}{5.5}
设 $X$ 为\Term{quasi-compact}{拟紧}\Term{quasi-separated}{拟分离}\Term{scheme}{概形}，$\ell\in\mathbb P$；$H$ 是定义在 $X$ 上\Term{p-torsion-sheaf}{$\ell$-层}\Term{category}{范畴}上的、取值于 $(\mathrm{Ab})$ 的半\Term{exact-functor}{正合函子}，并且与\Term{filtered}{滤过}\Term{inductive-limit}{归纳极限}相容。以下断言\Term{equivalence}{等价}：

(i) $H=0$。
\SourcePage{34}
(ii) 对每个\Term{finite-morphism}{有限态射} $f:Y\to X$，以及 $Y$ 的每个\Term{open-subset}{开集} $i:U\to Y$，其中 $U$ 在 $X$ 上为\Term{finite-presentation}{有限表现}，都有 $H(f_*i_!(\mathbb Z/\ell))=0$。若 $X$ \Term{noetherian}{诺特}，甚至只须取整的 $Y$。
\end{SGAresult}
\begin{proof}
由 2.7.2 和 2.5，只须对形如 $F=i_!G$ 的 $F$ 验证 $H(F)=0$；其中 $i:U\to X$ 是 $X$ 的一个在 $X$ 上\Term{finite-presentation}{有限表现}的子\Term{scheme}{概形}的包含，$G$ 为 $U$ 上\Term{locally-constant}{局部常值}且“\Term{isotrivial}{等平凡}”的\Term{p-torsion-sheaf}{$\ell$-层}。设 $U''\to U$ 为群为 $\mathfrak g$ 的\Term{principal-covering}{主覆盖}，使 $G$ 在 $U''$ 上成为\Term{constant-sheaf}{常层}。设 $A$ 为 $G$ 在 $U''$ 上取值的 $\ell$-群。群 $\mathfrak g$ \Term{action}{作用}于 $A$，而此\Term{action}{作用}决定 $G$ 的结构。令 $\mathfrak h$ 为 $\mathfrak g$ 的一个 Sylow $\ell$-\Term{subgroup}{子群}，$U'=U''/\mathfrak h$ 为对应 $\mathfrak h$ 的 $U$ 的\Term{covering}{覆盖}。$U'$ 在 $U$ 上的次数 $d$ 与 $\ell$ 互素。将 $f:U'\to U$ 延拓为一个\Term{finite-morphism}{有限态射} $g:X'\to X$，这样的延拓总存在（EGA IV 8.12.6）；令 $i':U'\to X'$ 为包含。应用 (5.4)，并利用 $d$ 与 $\ell$ 互素，归结为证明 $H(g_*i'_!f^*G)=0$。

众所周知，作为一个 $\ell$-群 $\mathfrak h$ 的模为单模的 $\ell$-挠\Term{abelian-group}{阿贝尔群}，只有赋予\Term{trivial-action}{平凡作用}的 $\mathbb Z/\ell$。因此，$A$ 作为 $\mathfrak h$-模存在一个\Term{filtered}{滤过}，其逐次商\Term{isomorphism}{同构}于赋予\Term{trivial-action}{平凡作用}的 $\mathbb Z/\ell$。此\Term{filtered}{滤过}给出 $f^*G$ 的一个相应\Term{filtered}{滤过}，所以为证明 $H(g_*i'_!f^*G)=0$，只须证明 $H(g_*i'_!\mathbb Z/r)=0$，结论即得。
\end{proof}
\begin{SGAresult}{命题}{5.6}
设 $X$ 为\Term{noetherian}{诺特}\Term{scheme}{概形}，$\ell\in\mathbb P$，$\varphi$ 为定义在 $X$ 上\Term{finite}{有限}整\Term{scheme}{概形}的\Term{set}{集合}上、取值于 $\mathbb N$ 的函数。设只要存在 $X$-\Term{morphism}{态射} $Y_1\to Y_2$，就有 $\varphi(Y_1)\le\varphi(Y_2)$；而若 $Y_1$ 为 $Y_2$ 的不同于 $Y_2$ 的闭子\Term{scheme}{概形}，则不等式严格。对每个 $X$ 上\Term{finite}{有限}的 $Y$，置
\[
\varphi(Y)=\sup_i\varphi(Y_i),
\]
其中 $Y_i$ 为 $Y$ 的各\Term{irreducible-component}{不可约分支}，赋予诱导的\Term{reduced}{既约}结构。则下列断言\Term{equivalence}{等价}：

(i) 对每个 $X$ 上\Term{finite}{有限}的 $Y$ 和 $Y$ 上的每个\Term{p-torsion-sheaf}{$\ell$-层} $F$，当 $q>\varphi(\operatorname{Supp}F)$ 时，有 $H^q(Y,F)=0$。

(ii) 对每个 $X$ 上\Term{finite}{有限}且整的 $Y$，当 $q>\varphi(Y)$ 时，有 $H^q(Y,\mathbb Z/\ell)=0$。
\end{SGAresult}
\begin{proof}
设 (ii) 成立。对每个 $H^q$ 可应用前一命题，并考虑到\Term{cohomology}{上同调}与\Term{inductive-limit}{归纳极限}以及\Term{finite-morphism}{有限态射}的\Term{direct-image}{直像}交换，因此验证 (i) 可归结到 $Y$ 在 $X$ 上\Term{finite}{有限}且整，而 $F$ 形如 $i_!(\mathbb Z/\ell)_U$ 的情形，其中 $i:U\to X$ 为\Term{nonempty}{非空}\Term{open-subset}{开集}。

有\Term{exact-sequence}{正合列}
\[
0\longrightarrow i_!(\mathbb Z/\ell)_U\longrightarrow(\mathbb Z/\ell)_Y\longrightarrow(\mathbb Z/\ell)_Z\longrightarrow0,
\]
其中 $Z=Y-U$，故 $Z\ne Y$。对 $\varphi$ 的值归纳，可设当 $q>\varphi(Y)-1$ 时 $H^q(Z,\mathbb Z/\ell)=0$；由\Term{cohomology}{上同调}\Term{exact-sequence}{正合列}，$q>\varphi(Y)$ 时 $H^q(Y,i_!(\mathbb Z/\ell))=0$。
\end{proof}
\SourcePage{35}
\begin{SGAresult}{推论}{5.7}
设 $X$ 为\Term{separably-closed-field}{可分闭域} $k$ 上的\Term{algebraic-curve}{代数曲线}，$k$ 的\Term{characteristic}{特征}不等于 $\ell$，$\ell\in\mathbb P$。则以\Term{p-torsion-sheaf}{$\ell$-挠层} $F$ 为值的 $X$ 的\Term{cohomology}{上同调}在\Term{dimension}{维数} $>2$ 时为零。若 $X$ \Term{affine}{仿射}，则\Term{cohomology}{上同调}在\Term{dimension}{维数} $>1$ 时为零。
\end{SGAresult}
\begin{proof}
对 $X$ 上\Term{finite}{有限}且整的 $Y$，置 $\varphi(Y)=2\dim Y$（若 $X$ 从而 $Y$ \Term{affine}{仿射}，则相应地置 $\varphi(Y)=\dim Y$）。由 7.6，只须证明当 $q>\varphi(Y)$ 时 $H^q(Y,\mathbb Z/\ell)=0$；当 $\dim Y=0$ 时这是平凡的。因此只须对整曲线 $Y$ 证明：当 $q>2$ 时（\Term{affine}{仿射}情形相应为 $q>1$）$H^q(Y,\mathbb Z/\ell)=0$。由于 $(\mathbb Z/\ell)_Y\simeq(\mu_\ell)_Y$，$q>2$ 的情形由 4.6 得到。

设 $X$ \Term{affine}{仿射}。由于\Term{radicial}{根性}\Term{group-extension}{扩张}不影响\Term{cohomology}{上同调}（第 VIII 讲 1.1），可归结到 $k$ \Term{algebraically-closed}{代数闭}的情形。设 $f:X'\to X$ 为 $X$ 的\Term{normalization}{正规化}，它为\Term{finite-morphism}{有限态射}，且除去某个\Term{finite}{有限}点集 $S$ 后为\Term{isomorphism}{同构}。$X$ 上有\Term{sheaf}{层}的\Term{exact-sequence}{正合列}
\[
0\longrightarrow(\mathbb Z/n)_X\longrightarrow f_*(\mathbb Z/n)_{X'}\longrightarrow\varepsilon\longrightarrow0,
\]
其中 $\varepsilon$ 集中在 $S$ 上，所以对 $q>0$，$H^q(X,\varepsilon)=0$。由第 VIII 讲 5.6，归结为证明：$q>1$ 时
\[
H^q(X,f_*(\mathbb Z/n)_{X'})\simeq H^q(X',\mathbb Z/n)=0,
\]
即归结到 $X$ \Term{normal}{正规的}情形。显然还可设 $X$ \Term{connected}{连通}。设 $i:X\to\overline X$ 为\Term{open-immersion}{开浸入}，其中 $\overline X$ \Term{complete}{完备}且\Term{nonsingular}{非奇异}。由于 $X$ \Term{affine}{仿射}，$\overline X$ 中至少缺去一个点，因此立即看出，\Term{morphism}{态射}
\[
A(k)\longrightarrow\operatorname{Pic}X
\]
\Term{surjectivity}{满射}，其中 $A$ 为 $\overline X$ 的 Jacobian（参看 (4.8)）。由于 $A(k)$ 对 $\ell\ne\operatorname{char}k$ 可除，$\operatorname{Pic}X$ 也如此，故由 4.6，$H^2(X,\mu_\ell)=0$。证毕。
\end{proof}
\textbf{注记 5.8。} 为供以后引用，指出 5.5 的论证实际上证明了下述结果：设 $X$ 为\Term{quasi-compact}{拟紧}\Term{quasi-separated}{拟分离}\Term{scheme}{概形}，$\ell$ 为素数，$\mathcal C$ 是\Term{constructible}{可构造}的\Term{p-torsion-sheaf}{$\ell$-挠}\Term{abelian-sheaf}{阿贝尔层}\Term{category}{范畴}的一个\Term{strictly-full-subcategory}{严格满子范畴}，对\Term{direct-summand}{直和因子}及\Term{group-extension}{扩张}封闭。设对每个\Term{finite-morphism}{有限态射} $f:Y\to X$，以及 $Y$ 的每个在 $X$ 上为\Term{finite-presentation}{有限表现}的\Term{open-subset}{开集} $i:U\to Y$，有 $f_*(i_!(\mathbb Z/\ell\mathbb Z)_U)\in\mathcal C$。则 $\mathcal C$ 包含所有\Term{constructible}{可构造}\Term{p-torsion-sheaf}{$\ell$-挠层}。

\SGAsection{}{参考文献}
\begin{enumerate}
\item[{[1]}] Gabriel, P.《\Term{abelian-category}{阿贝尔范畴}》（\textit{Des Catégories Abéliennes}），博士论文，巴黎，1962。
\item[{[2]}] Ogg, A.《\Term{function-field}{函数域}上的\Term{abelian-variety}{阿贝尔簇}的\Term{cohomology}{上同调}》（\textit{Cohomology of abelian varieties over functions fields}），\textit{Annals of Math.}，第 76 卷，第 2 期，1962。
\item[{[3]}] Deuring, M.《椭圆\Term{function-field}{函数域}的\Term{multiplier}{乘子}环的类型》（\textit{Die Typen der Multiplikatorenringe elliptischer Funktionenkörper}），\textit{Hamb. Abh. Bol.}，第 13 卷，1940，第 197 页。
\SourcePage{36}
\item[{[4]}] Grothendieck, A.《\Term{algebra}{代数}几何基础》（\textit{Fondements de la Géométrie Algébrique}），Bourbaki 讨论班 1957--1962 年选文，Secrétariat mathématique，11, rue Pierre Curie, Paris 5ème。（简称 FGA。）
\item[{[5]}] Serre, J.-P.《Galois \Term{cohomology}{上同调}》（\textit{Cohomologie Galoisienne}），Lecture Notes 第 5 号，Springer。（简称 CG。）
\item[{[6]}] Douady, A.《紧\Term{totally-disconnected}{全不连通}群的\Term{cohomology}{上同调}》（\textit{Cohomologie des groupes compacts totalement discontinus}），Bourbaki 讨论班，第 189 讲。
\end{enumerate}
\SourcePage{37}
% 原页空白。
